\chapter{Image de jauge et composition vers Yang--Mills}

\section{Objet source et séparation obligatoire des générateurs}

L’objet déjà engendré que reçoit ce chapitre est la structure de jauge complète
$\mathfrak G_{\rm gau}$ ; il ne constitue ni une donnée externe ni une hypothèse terminale.
À chaque niveau $m$, il conserve le réseau
$\Lambda_m=a_m\mathbb Z^4$, $a_{m+1}=a_m/9$, la mesure $\mu_m$, les
opérateurs d’entrelacement $J_m$, quatre réflexions $\Theta_{\nu,m}$, l’holonomie, la couleur,
la représentation, la retenue de fusion, la route, la mémoire, le terminal et le lecteur de courbure.
Ses deux dynamiques sont typées par des symboles distincts avant toute limite.

Le semi-groupe d’une seule fibre de mémoire est
\begin{equation}\label{eq:pdf2-ym-fibra-semigrupo}
 K^{\rm fib}_{m,t}
 =P_{\Omega_m}+e^{-3\kappa_{\rm av}t}(I-P_{\Omega_m})
 =e^{-tD_m^{\rm fib}},
 \qquad
 D_m^{\rm fib}:=3\kappa_{\rm av}(I-P_{\Omega_m}).
\end{equation}
Son spectre est contenu dans $\{0,3\kappa_{\rm av}\}$. Le générateur de la carte produit complète est, en revanche,
\begin{equation}\label{eq:pdf2-ym-generador-producto}
 D_m^{\rm prod}:=3\kappa_{\rm av}
 \sum_{c\in\mathcal I_m}(I-E_{m,c}),
 \qquad
 K^{\rm prod}_{m,t}:=e^{-tD_m^{\rm prod}},
\end{equation}
où les espérances orthogonales $E_{m,c}$ commutent. S’il y a deux ou plusieurs
coordonnées non constantes, $D_m^{\rm prod}$ possède des niveaux
$6\kappa_{\rm av},9\kappa_{\rm av},\ldots$ ; c’est pourquoi
\eqref{eq:pdf2-ym-fibra-semigrupo} ne peut être déclaré égal à
\eqref{eq:pdf2-ym-generador-producto}. Le premier est la fibre simple ; le second,
le produit de recouvrement.

\section{Théorème fini de Hoeffding, vide et écart spectral atteint}

Dans la carte produit construite dans \eqref{eq:pdf2-ym-owner-product-factorization}, les $E_{m,c}$ sont des espérances conditionnelles orthogonales et commutantes, et la preuve de \eqref{eq:pdf2-ym-owner-vacuum} établit
\begin{equation}\label{eq:pdf2-ym-interseccion-vacio}
 \bigcap_{c\in\mathcal I_m}\operatorname{Ran}E_{m,c}
 =\mathbb C\Omega_m.
\end{equation}
La décomposition de Hoeffding est la somme orthogonale
\[
 \mathcal H_m=\bigoplus_{S\subseteq\mathcal I_m}\mathcal H_{m,S},
 \qquad u=\sum_Su_S,
\]
où $E_{m,c}u_S=0$ si $c\in S$ et $E_{m,c}u_S=u_S$ si $c\notin S$. Alors
\begin{equation}\label{eq:pdf2-ym-hoeffding}
 D_m^{\rm prod}u_S=3\kappa_{\rm av}|S|u_S,
 \qquad
 \langle u,D_m^{\rm prod}u\rangle
 =3\kappa_{\rm av}\sum_S|S|\|u_S\|^2.
\end{equation}
La composante $S=\varnothing$ est exactement le vide de
\eqref{eq:pdf2-ym-interseccion-vacio}. Par conséquent
\begin{equation}\label{eq:pdf2-ym-cota-finita}
 D_m^{\rm prod}\succeq
 3\kappa_{\rm av}(I-P_{\Omega_m}),
 \qquad
 \|K^{\rm prod}_{m,t}(I-P_{\Omega_m})\|
 \le e^{-3\kappa_{\rm av}t}.
\end{equation}

Dans une coordonnée d’incidence
$q=(q_1,\ldots,q_6)\in\mathbb F_3^6$, le témoin
\begin{equation}\label{eq:pdf2-ym-wc}
 W_c(q):=\mathbf1_{\{q_1+\cdots+q_6=0\}}-\frac13
\end{equation}
satisfait
\begin{equation}\label{eq:pdf2-ym-wc-momentos}
 \mathbb EW_c=0,\qquad \|W_c\|^2=\frac29,\qquad
 D_m^{\rm prod}W_c=3\kappa_{\rm av}W_c.
\end{equation}
Ainsi, à chaque niveau fini, le seuil de \eqref{eq:pdf2-ym-cota-finita} est atteint. Ce calcul fixe le spectre du générateur produit fini ; la mesure quadridimensionnelle continue de Yang--Mills est construite matériellement dans \eqref{eq:pdf2-ym-gibbs-measure}--\eqref{eq:pdf2-ym-g-dynamics}.

\section{Compression de jauge physique}

Soit $G$ compact, connexe et simple. Le groupe qui agit au niveau $m$ est
\[
 \mathcal G_m^{\rm full}
 =G^{V(\Lambda_m)}\times
  \prod_{c\in\mathcal C_m}\mathbb F_3^4,
\]
où le second facteur agit sur les six incidences au moyen d’une matrice
qui doit rester visible. Si les lignes sont indexées par
$(01,02,03,12,13,23)$ et les colonnes par $0,1,2,3$, on fixe
\begin{equation}\label{eq:pdf2-ym-incidence-matrix}
 B^+_{(ij),k}:=\delta_{ik}+\delta_{jk}\in\mathbb F_3,
 \qquad
 (B^+x)_{ij}=x_i+x_j.
\end{equation}
C’est l’incidence non orientée du collier ; elle ne se confond pas avec le
cobord signé d’une arête. L’action finie est
\begin{equation}\label{eq:pdf2-ym-incidence-action}
 q_c\longmapsto q_c+B^+x_c,
 \qquad x_c\in\mathbb F_3^4.
\end{equation}
La projection de Haar
\begin{equation}\label{eq:pdf2-ym-proyeccion-fisica}
 \mathbb P_m^{\rm phys}F
 :=\int_{\mathcal G_m^{\rm full}}F(g^{-1}\!\cdot\,)\,dg,
 \qquad
 \mathcal H_m^{\rm phys}:=\operatorname{Ran}\mathbb P_m^{\rm phys}
\end{equation}
est distincte de la marginale qui oublie l'incidence. L'invariance des espérances $E_{m,c}$ donne
\begin{equation}\label{eq:pdf2-ym-reduccion}
 [\mathbb P_m^{\rm phys},D_m^{\rm prod}]=0,
 \qquad
 [\mathbb P_m^{\rm phys},K_{m,t}^{\rm prod}]=0.
\end{equation}
De plus, chaque variable de $x_c\in\mathbb F_3^4$ apparaît trois fois dans la somme des
six incidences, de sorte que
\begin{equation}\label{eq:pdf2-ym-incidence-conservation}
 \sum_{i<j}(B^+x_c)_{ij}
 =3\sum_{i=0}^3x_{c,i}=0
 \quad\text{dans }\mathbb F_3.
\end{equation}
Par conséquent, $W_c$ est invariant de
jauge et appartient à $\mathcal H_m^{\rm phys}$. L’écart fini n’est pas un
artefact d’une dilatation non physique.

\section{Positivité par réflexion à chaque niveau}

Pour un noyau réversible commun, définissons la mesure à huit faces
\begin{equation}\label{eq:pdf2-ym-ocho-caras}
 \Omega_m^{(8)}(dy_1\cdots dy_8)
 :=\int\mu_m(dz)\prod_{r=1}^{8}K^{\rm prod}_{m,a_m/2}(z,dy_r).
\end{equation}
Si $\Theta_{\nu,m}$ échange les deux moitiés par rapport à la direction $\nu$, la propriété de Markov et la réversibilité factorisent, pour tout observable cylindrique $F$ à support dans le demi-espace positif,
\begin{equation}\label{eq:pdf2-ym-os-finita}
 \int\overline{F(\Theta_{\nu,m}y)}F(y)\,d\Omega_m^{(8)}
 =\|\mathcal B_{m,\nu}F\|_2^2\ge0,
 \qquad \nu=1,\ldots,4.
\end{equation}
La marginale de faces opposées récupère le même
$K^{\rm prod}_{m,a_m}$. Ainsi, les quatre formes OS finies sont
la conséquence d’une seule mesure et d’un seul semi-groupe, non de quatre processus
indépendants.

\section{Propriétaire final : l'état de jauge complet, non sa marginale amputée}

Le propriétaire en vigueur ne prend pas pour espace physique la marginale qui ne conserve que
les liens. Son domaine est $\mathfrak G_{\rm gau}$, avec lien, repère,
représentation, incidence, couleur, route, orientation,
retenue, mémoire, survie et terminal. La projection vers les liens est une application
d’oubli ultérieure. Juger le témoin physique après avoir appliqué cet oubli remplace
l’objet et ne transfère pas la conclusion.

\subsection{Mesure projective et publication des liaisons}

Pour un écran fini $\Gamma_m=(V_m,E_m)$, on part de l'espace relevé
\begin{equation}\label{eq:pdf2-ym-owner-lifted-space}
 \widetilde X_m:=G^{V_m}\times G^{E_m},\qquad
 d\widetilde\mu_m(h,z)
 =\prod_{v\in V_m}dh_v\prod_{e\in E_m}k_{\tau_{m,e}}(z_e)\,dz_e,
\end{equation}
et chaque lien est publié au moyen de
\begin{equation}\label{eq:pdf2-ym-owner-link-publication}
 U_e=h_{s(e)}z_eh_{t(e)}^{-1}.
\end{equation}
Une arête grossière $e$ se raffine en la chaîne ordonnée $e_1\cdots e_9$. Posons \(\boldsymbol\tau=(\tau_1,\ldots,\tau_9)\), avec \(\tau_j=\tau_{m+1,e_j}>0\) et \(\tau=\tau_{m,e}=\sum_{j=1}^9\tau_j\). Soit \(d\lambda_z^{(9)}\) la désintégration de Haar sur la fibre du produit, caractérisée par
\[
 \int_{G^9}F(z_1,\ldots,z_9)\,dz_1\cdots dz_9
 =\int_G\int_{z_1\cdots z_9=z}
   F\,d\lambda_z^{(9)}\,dz.
\]
Le pont correct pour des temps non nécessairement égaux est
\begin{equation}\label{eq:pdf2-ym-owner-heat-bridge}
 dB_{z,\boldsymbol\tau}^{(9)}(z_1,\ldots,z_9)
 =\frac{\prod_{j=1}^9k_{\tau_j}(z_j)}{k_\tau(z)}
   \,d\lambda_z^{(9)},
 \qquad z_1\cdots z_9=z,
\end{equation}
qui est une probabilité car
\(k_{\tau_1}*\cdots*k_{\tau_9}=k_\tau\). Dans le raffinement nonadique uniforme, on
récupère \(\tau_j=\tau/9\), mais la formule ne présuppose pas cette égalité.

Pour rendre l'innovation indépendante sans effacer le produit grossier, on fixe, pour chaque $(e,z)$, un transport borélien
\begin{equation}\label{eq:pdf2-ym-owner-bridge-triangularization}
 \chi_{m,e,z}:\bigl((0,1)^{\mathbb N},du^{\otimes\mathbb N}\bigr)
 \longrightarrow
 \bigl(\{(z_j):z_1\cdots z_9=z\},B_{z,\boldsymbol\tau}^{(9)}\bigr)
\end{equation}
qui est un isomorphisme modulo des ensembles nuls. Il existe car toutes deux sont des réalisations
standard non atomiques ; un choix est fixé une fois pour toutes et transporté sous les
réétiquetages des arêtes. La variable grossière $z$ et l’innovation $u^{\rm det}_{m+1,e}$
sont alors des facteurs indépendants dans la carte triangulaire. Les nouveaux repères sont
intégrés avec Haar. La route, la retenue, l’orientation et la mémoire sont transportées par les
morphismes $\operatorname{Ext}_{\rm gau}$ de la donnée de jauge elle-même.

La projection grossière multiplie les neuf accroissements et oublie exclusivement les
innovations nouvelles. Avec cette définition, on obtient
\begin{equation}\label{eq:pdf2-ym-owner-projectivity}
 (\widehat\pi_{m+1,m})_*\widehat\mu_{m+1}^{\rm ov}
 =\widehat\mu_m^{\rm ov},\qquad
 \widehat J_mF=F\circ\widehat\pi_{m+1,m}.
\end{equation}
La même arête apparaît une seule fois dans la mesure ; les six faces incidentes sont six lecteurs orientés de cette arête. Les chemins qui se recouvrent conservent l'accroissement partagé et ne sont pas remplacés par des copies indépendantes.

\subsection{La fibre d’incidence est une coordonnée physique de l’état complet}

La queue d’incidence incluse dans $\mathfrak G_{\rm gau}$ se décompose sans perdre
d’information :
\begin{equation}\label{eq:pdf2-ym-owner-seed-split}
 \Sigma:(\mathbb F_3^2)^{\mathbb N}\longrightarrow
 (\mathbb F_3^2)^{\mathbb N}\times\mathbb F_3^6,
 \qquad \Sigma_*\sigma=\sigma\otimes u_6.
\end{equation}
Le premier facteur continue d’alimenter les cartes d’innovation et le second conserve
les six incidences $q_c=(q_{01},q_{02},q_{03},q_{12},q_{13},q_{23})$ du collier
$c$. Au lieu de laisser implicite quelles coordonnées le générateur gouverne, on fixe une
carte triangulaire récursive
\begin{equation}\label{eq:pdf2-ym-owner-product-factorization}
 \widehat X_m
 =\widehat X_{m_0}^{\rm seed}
  \times\prod_{\ell=m_0+1}^{m}
       \mathcal U_{\ell/\ell-1}^{\rm det},
 \qquad
 \widehat\mu_m^{\rm ov}
 =\widehat\mu_{m_0}^{\rm seed}
  \otimes\bigotimes_{\ell=m_0+1}^{m}
       \upsilon_{\ell/\ell-1}^{\rm det}.
\end{equation}
Le facteur germe énumère, une fois, les repères et les liaisons grossières, les incidences $q_c$ et leurs cartes de marquage. Chaque $\mathcal U_{\ell/\ell-1}^{\rm det}$ énumère, une fois, les nouveaux repères, les ponts \eqref{eq:pdf2-ym-owner-bridge-triangularization}, les nouvelles incidences et les sous-queues de marquage de ce raffinement. La factorisation est interne à $\mathfrak G_{\rm gau}$ et n'ajoute pas d'autre donnée de départ.

On enregistre l’index exhaustif, disjoint et dénombrable
\begin{equation}\label{eq:pdf2-ym-owner-exhaustive-index}
 \mathcal I_m
 :=\mathcal I_{m_0}^{\rm frame}
 \sqcup\mathcal I_{m_0}^{\rm link}
 \sqcup\mathcal I_{m_0}^{q}
 \sqcup\mathcal I_{m_0}^{\rm mark}
 \sqcup\bigsqcup_{\ell=m_0+1}^{m}
 \bigl(
   \mathcal I_{\ell/\ell-1}^{\rm frame}
   \sqcup\mathcal I_{\ell/\ell-1}^{\rm bridge}
   \sqcup\mathcal I_{\ell/\ell-1}^{q}
   \sqcup\mathcal I_{\ell/\ell-1}^{\rm mark}
 \bigr).
\end{equation}
Chaque facteur indépendant de \eqref{eq:pdf2-ym-owner-product-factorization} apparaît
exactement une fois ; une sous-queue infinie est énumérée par ses cylindres scalaires. Aucun
facteur généalogique ne reste hors de $\mathcal I_m$.
En particulier, les abréviations de la récurrence signifient littéralement
\begin{equation}\label{eq:pdf2-ym-owner-index-abbreviations}
 \begin{aligned}
 \mathcal I_{m_0}^{\rm seed}
 &:={\mathcal I}_{m_0}^{\rm frame}\sqcup
    {\mathcal I}_{m_0}^{\rm link}\sqcup
    {\mathcal I}_{m_0}^{q}\sqcup
    {\mathcal I}_{m_0}^{\rm mark},\\
 \mathcal I_{\ell/\ell-1}^{\rm det}
 &:={\mathcal I}_{\ell/\ell-1}^{\rm frame}\sqcup
    {\mathcal I}_{\ell/\ell-1}^{\rm bridge}\sqcup
    {\mathcal I}_{\ell/\ell-1}^{q}\sqcup
    {\mathcal I}_{\ell/\ell-1}^{\rm mark}.
 \end{aligned}
\end{equation}

La marginale des liens oublie $q_c$ ; la sous-algèbre physique, en revanche, s’obtient
en moyennant l’action de jauge complète
\begin{equation}\label{eq:pdf2-ym-full-gauge-group}
 \mathcal G_m^{\rm full}
 :=G^{V(\Lambda_m)}\times
 \prod_{c\in\mathcal C_m}\mathbb F_3^4,
 \qquad q_c\longmapsto q_c+B^+x_c.
\end{equation}
La projection physique est donc
\begin{equation}\label{eq:pdf2-ym-full-physical-projection}
 \mathbb P_m^{\rm phys}F
 :=\int_{\mathcal G_m^{\rm full}}F(g^{-1}\!\cdot x)\,dg,
 \qquad
 \widehat{\mathscr H}_m^{\rm phys}
 :=\operatorname{Ran}\mathbb P_m^{\rm phys}.
\end{equation}
Elle ne coïncide pas avec l’application d’oubli
$\widehat X_m\to X_m^{\rm link}$. Cette distinction est le garde-fou de typage
qui empêche d’expulser $W_c$ du secteur physique.

\section{Générateur propriétaire et circuit local}

La carte triangulaire complète fournit une application unitaire
\begin{equation}\label{eq:pdf2-ym-owner-triangular-unitary}
 \widehat\Gamma_m:
 L^2(\widehat X_{m+1},\widehat\mu_{m+1}^{\rm ov})
 \longrightarrow
 L^2(\widehat X_m,\widehat\mu_m^{\rm ov})
 \widehat\otimes\mathcal K_{m+1}^{\rm det},
 \qquad
 \widehat\Gamma_m\widehat J_mF=F\otimes\mathbf1_{\rm det}.
\end{equation}
Pour $\iota\in\mathcal I_m$, soit $E_{m,\iota}$ l'espérance qui intègre exactement le facteur $\iota$ de \eqref{eq:pdf2-ym-owner-product-factorization} et laisse tous les autres fixes. Ce sont des projections de Markov orthogonales, elles conservent l'unité et commutent. Sur le cœur dense $\mathscr C_m$ des cylindres qui dépendent d'un nombre fini de coordonnées, on définit la forme
\begin{equation}\label{eq:pdf2-ym-owner-full-form}
 \widehat{\mathcal E}_m[F]
 :=3\kappa_{\rm av}\sum_{\iota\in\mathcal I_m}
   \|(I-E_{m,\iota})F\|_2^2.
\end{equation}
La somme est finie pour chaque cylindre. Sa fermeture est une forme de Dirichlet et
détermine le générateur auto-adjoint markovien
$\widehat H_m^{\rm ov}$. À un niveau germe, cela s’écrit, au sens des
formes,
\begin{equation}\label{eq:pdf2-ym-owner-base-generator}
 \widehat H_{m_0}^{\rm ov}
 =3\kappa_{\rm av}
  \sum_{\iota\in\mathcal I_{m_0}^{\rm seed}}(I-E_{m_0,\iota}),
\end{equation}
et chaque raffinement ajoute exactement les espérances de ses nouvelles coordonnées :
\begin{equation}\label{eq:pdf2-ym-owner-refined-generator}
 \begin{aligned}
 \widehat H_{m+1}^{\rm ov}
 &=\widehat\Gamma_m^*
   (\widehat H_m^{\rm ov}\otimes I
    +I\otimes\widehat H_{m+1}^{\rm det})\widehat\Gamma_m,\\
 \widehat H_{m+1}^{\rm det}
 &=3\kappa_{\rm av}
   \sum_{\iota\in\mathcal I_{m+1/m}^{\rm det}}
   (I-E_{m+1,\iota}^{\rm det}).
 \end{aligned}
\end{equation}
La partition des couleurs n’est pas laissée nominale. Si une cellule est indexée par
$n=(n_1,n_2,n_3,n_4)\in\mathbb Z^4$, on fixe
\begin{equation}\label{eq:pdf2-ym-owner-729-coloring}
 \chi_m(n):=
 \bigl(n_1\bmod9,\ n_2\bmod9,\ n_3+3n_4\bmod9\bigr)
 \in(\mathbb Z/9\mathbb Z)^3,
 \qquad |(\mathbb Z/9\mathbb Z)^3|=9\cdot81.
\end{equation}
Si deux colliers fermés de rayon cellulaire un se coupent, leur différence
$\delta$ vérifie $|\delta_j|\le2$. L’égalité des couleurs impose
$\delta_1=\delta_2=0$ et
$\delta_3+3\delta_4\equiv0\pmod9$. Comme le dernier entier appartient à
$[-8,8]$, il vaut zéro ; les bornes $|\delta_3|,|\delta_4|\le2$ donnent alors
$\delta_3=\delta_4=0$. Par conséquent, des cellules distinctes de même couleur ont
des colliers disjoints. On définit matériellement
\begin{equation}\label{eq:pdf2-ym-owner-729-blocks}
 \mathcal B_{m,r}:=\{\iota\in\mathcal I_m:
   \operatorname{cell}(\iota)=n, \chi_m(n)=r\},
 \qquad r\in(\mathbb Z/9\mathbb Z)^3.
\end{equation}
Le transport d’un écran par $g\in E(4)$ transporte aussi cette
partition ; il n’exige pas qu’une rotation préserve les coordonnées de l’étiquette.
Ainsi, $\{\mathcal B_{m,r}\}_{r=1}^{9\cdot81}$ regroupe des supports spatiaux
disjoints ;
l’indice interne de chaque bloc énumère ses facteurs de
\eqref{eq:pdf2-ym-owner-exhaustive-index}. Pour un cylindre $F$, le produit ne
contient que les blocs dont dépend $F$ ; la fermeture définit le produit fort
pour toute la forme. Le générateur d’un bloc conserve les multiplicités de
Hoeffding ; il n’est pas remplacé par un seul projecteur qui les effacerait. Par
commutativité dans chaque carte,
\begin{equation}\label{eq:pdf2-ym-owner-local-circuit}
\begin{aligned}
 H_{m,B}&:=3\kappa_{\rm av}\sum_{\iota\in B}(I-E_{m,\iota}),\\
 K_{m,B}&:=e^{-a_mH_{m,B}}
 =\prod_{\iota\in B}e^{-3\kappa_{\rm av}a_m(I-E_{m,\iota})},\\
 K_m^{\rm loc}&:=\mathop{\prod_{r,B}}^{\rm strong}K_{m,B}
 =e^{-a_m\widehat H_m^{\rm ov}}.
\end{aligned}
\end{equation}
Le diamètre physique de chaque bloc reste uniformément borné. C'est de là que provient la localité : le générateur n'introduit pas de mise à jour instantanée sur tout le réseau.
La partition est compatible avec Bianchi, et pas seulement avec la disjonction géométrique. Pour une cellule cubique $c$, on fixe le produit orienté, avec tous les facteurs transportés au même sommet,
\begin{equation}\label{eq:pdf2-ym-owner-bianchi-block-product}
 \mathscr B_{m,c}(x):=
 \prod_{p\subset\partial c}^{\longrightarrow}
 T_p\operatorname{Hol}_p(x)T_p^{-1}.
\end{equation}
Chaque arête interne de la frontière apparaît deux fois avec des orientations contraires ;
par conséquent, $\mathscr B_{m,c}=e$ point par point. Un bloc
$\mathcal B_{m,r}$ contient le collier complet de toute cellule qu’il met à jour : s’il touche
une face de $c$, il met simultanément à jour la paire inverse qui partage cette arête.
L’annulation précédente reste littérale après chaque facteur du circuit :
\begin{equation}\label{eq:pdf2-ym-owner-bianchi-block-invariance}
 K_{m,B}\mathbf1_{\{\mathscr B_{m,c}=e\}}
 =\mathbf1_{\{\mathscr B_{m,c}=e\}},
 \qquad
 K_m^{\rm loc}\mathbf1_{\{\mathscr B_{m,c}=e\}}
 =\mathbf1_{\{\mathscr B_{m,c}=e\}}.
\end{equation}
Ainsi, le calendrier de $9\cdot81$ couches préserve Bianchi à chaque étape et dans le produit fort, et pas seulement à la fin d'un balayage.

L'action de jauge permute ou translate les coordonnées intégrées par chaque $E_{m,\iota}$ et préserve leurs mesures. Par conséquent
\begin{equation}\label{eq:pdf2-ym-owner-physical-reduction}
 U_m(g)\widehat H_m^{\rm ov}=\widehat H_m^{\rm ov}U_m(g)
 \quad(g\in\mathcal G_m^{\rm full}),
 \qquad
 [\mathbb P_m^{\rm phys},\widehat H_m^{\rm ov}]=0,
 \qquad
 \mathbb P_{m+1}^{\rm phys}\widehat J_m
 =\widehat J_m\mathbb P_m^{\rm phys}.
\end{equation}
Le semi-groupe complet et son noyau de Markov sont notés
\begin{equation}\label{eq:pdf2-ym-owner-full-kernel}
 \widehat K_{m,t}^{\rm ov}:=e^{-t\widehat H_m^{\rm ov}},
 \qquad
\widehat K_{m,t}^{\rm ov}(x,dy)
 \quad(x,y\in\widehat X_m).
\end{equation}
La première identité de \eqref{eq:pdf2-ym-owner-physical-reduction} implique
\begin{equation}\label{eq:pdf2-ym-owner-full-kernel-equivariance}
 \widehat K_{m,t}^{\rm ov}(gx,gA)
 =\widehat K_{m,t}^{\rm ov}(x,A)
 \quad(g\in\mathcal G_m^{\rm full}).
\end{equation}
Pour typer la réduction sans utiliser un noyau restreint sur le mauvais espace,
soit $q_m:\widehat X_m\to Y_m:=\widehat X_m/\mathcal G_m^{\rm full}$ l’application des
orbites et $\nu_m=(q_m)_*\widehat\mu_m^{\rm ov}$. Le tiré en arrière
\begin{equation}\label{eq:pdf2-ym-owner-quotient-unitary}
 \mathcal W_m:L^2(Y_m,\nu_m)\longrightarrow
 \widehat{\mathscr H}_m^{\rm phys},
 \qquad \mathcal W_m\phi=\phi\circ q_m
\end{equation}
est unitaire. L'équivariance du noyau complet permet de définir
\begin{equation}\label{eq:pdf2-ym-owner-physical-kernel}
 K_{m,t}^{\rm phys}(q_mx,A)
 :=\widehat K_{m,t}^{\rm ov}(x,q_m^{-1}A),
 \qquad A\subseteq Y_m,
\end{equation}
indépendamment du représentant $x$. L’objet physique est alors typé par
\begin{equation}\label{eq:pdf2-ym-owner-D}
 D_m^{\rm YM}
 :=\left.\widehat H_m^{\rm ov}\right|_{\widehat{\mathscr H}_m^{\rm phys}},
 \qquad
 e^{-tD_m^{\rm YM}}
 =\mathcal W_mK_{m,t}^{\rm phys}\mathcal W_m^{-1}.
\end{equation}
C'est le générateur produit complet ; ce n'est pas le générateur d'une seule fibre de \eqref{eq:pdf2-ym-fibra-semigrupo}.

\section{Action de Yang--Mills, mesure dépendant de $g$ et transport exact}

Jusqu’ici, la carte cinématique de référence a été écrite. Nous la noterons
$\widehat\mu_m^0:=\widehat\mu_m^{\rm ov}$ et
$\widehat H_m^0:=\widehat H_m^{\rm ov}$, avec
$\widehat K_{m,t}^0:=e^{-t\widehat H_m^0}=\widehat K_{m,t}^{\rm ov}$.
Pour transformer le lecteur structurel
$\mathcal P_m^{F^2}$ en une loi, et non en une comparaison ultérieure, on reporte
d’abord sa racine positive sur la factorisation triangulaire. La polarisation de
$\mathcal P_m^{F^2}$ dans un collier $b$ est le Gram positif
\[
 Q_{m,b}(u,v):=\frac14\bigl(
 \mathcal P_{m,b}^{F^2}(u+v)-\mathcal P_{m,b}^{F^2}(u-v)
 \bigr).
\]
On quotiente l’espace fini des lecteurs du collier par $\ker Q_{m,b}$, on le complète
avec $Q_{m,b}$ et on note $\mathcal E_{m,b}^{F}$ le Hilbert résultant. Si
$\zeta_{m,b}(x_b)$ est le vecteur des lecteurs centrés de la coordonnée triangulaire
$x_b$, on fixe matériellement
\begin{equation}\label{eq:pdf2-ym-action-root}
 \xi_{m,b}(x_b):=Q_{m,b}^{1/2}\zeta_{m,b}(x_b)
 \in\mathcal E_{m,b}^{F},
 \qquad
 p_{m,b}(x_b):=\|\xi_{m,b}(x_b)\|^2.
\end{equation}
L'indice $b$ parcourt les colliers de Gram à support minimal. Un collier est affecté à la première couleur de \eqref{eq:pdf2-ym-owner-729-coloring} qui contient sa cellule de base ; ainsi, chaque terme apparaît une fois et les colliers d'une même couleur sont disjoints. Plus précisément, la diagonalisation positive regroupe les facteurs triangulaires en une partition disjointe
\begin{equation}\label{eq:pdf2-ym-action-factor-partition}
 \mathcal I_m^{F}
 =\bigsqcup_{r\in(\mathbb Z/9\mathbb Z)^3}
   \bigsqcup_{b\in\mathcal B_{m,r}^{F}}\mathcal I_{m,b},
 \qquad x_b:=(x_\iota)_{\iota\in\mathcal I_{m,b}},
 \qquad
 \mathcal B_{m,r}^{F}:=\{b:Q_{m,b}\ne0,\ \chi_m(b)=r\}.
\end{equation}
Aucune coordonnée triangulaire n’appartient à deux $\mathcal I_{m,b}$ ; si un lecteur
géométrique touche deux colliers, son mode de Gram est assigné au premier et son complément
orthogonal au second. C’est la décomposition spectrale finie de la forme positive,
non une duplication de la variable physique.
La racine de la somme orthogonale, calculée couleur par couleur, donne l’identité, non une
approximation,
\begin{equation}\label{eq:pdf2-ym-action-reader-sum}
 \mathcal P_m^{F^2}(x)
 =\sum_{r\in(\mathbb Z/9\mathbb Z)^3}
   \sum_{b\in\mathcal B_{m,r}^{F}}p_{m,b}(x_b),
 \qquad
 S_{m,g}(x):=\frac1{4g^2}\mathcal P_m^{F^2}(x),
 \quad g\in(0,\infty).
\end{equation}
Ici, $x_b:\widehat X_m\to X_{m,b}$ est la projection des coordonnées du
collier, de sorte que le domaine, l’action et le codomaine de chaque terme sont fixés.
Le lecteur est invariant de jauge parce que $Q_{m,b}$ utilise le produit invariant de
$\mathfrak g$ ; il est invariant par réflexion parce qu’une réflexion ne fait que permuter les
six paires orientées ; et il est naturel par réétiquetage euclidien.
La partition \eqref{eq:pdf2-ym-action-factor-partition} et la carte produit donnent
explicitement
\begin{equation}\label{eq:pdf2-ym-action-factor-measures}
 \widehat\mu_m^0=\lambda_m^{\rm inert}\otimes
   \bigotimes_{r,b}\lambda_{m,b}^0,
 \qquad
 d\lambda_{m,b}^g
 :=(Z_{m,b}^g)^{-1}e^{-p_{m,b}/(4g^2)}d\lambda_{m,b}^0,
 \qquad
 \widehat\mu_{m,g}^{\rm YM}=\lambda_m^{\rm inert}\otimes
   \bigotimes_{r,b}\lambda_{m,b}^g.
\end{equation}
Le facteur inerte contient exactement les coordonnées sur lesquelles le lecteur est nul ; il n'est pas éliminé de la théorie.

La séparation grossier--détail de \eqref{eq:pdf2-ym-owner-product-factorization} est orthogonale pour la même racine : la polarisation directe donne $Q_{m+1}=Q_m\oplus Q_{m+1/m}^{\rm det}$ et $\zeta_{m+1}=\zeta_m\oplus\zeta_{m+1/m}^{\rm det}$. Si $x_{m+1}=\widehat\Gamma_m^{-1}(x_m,u)$, alors
\begin{equation}\label{eq:pdf2-ym-action-cocycle}
 \xi_{m+1}(x_{m+1})
 =\xi_m(x_m)\oplus\xi_{m+1/m}^{\rm det}(u),
 \qquad
 S_{m+1,g}\circ\widehat\Gamma_m^{-1}
 =S_{m,g}+S_{m+1/m,g}^{\rm det}.
\end{equation}
C'est la forme précise sous laquelle retenue et mémoire empêchent de compter deux fois une face ancestrale : la partie grossière n'est pas à nouveau additionnée dans le détail. En particulier, $Z_{m+1,g}=Z_{m,g}Z_{m+1/m,g}^{\rm det}$ pour les normalisateurs et l'on obtient la famille de probabilités véritablement dépendante du couplage
\begin{equation}\label{eq:pdf2-ym-gibbs-measure}
 d\widehat\mu_{m,g}^{\rm YM}(x)
 :=Z_{m,g}^{-1}e^{-S_{m,g}(x)}d\widehat\mu_m^0(x),
 \qquad
 (\widehat\pi_{m+1,m})_*\widehat\mu_{m+1,g}^{\rm YM}
 =\widehat\mu_{m,g}^{\rm YM}.
\end{equation}
Pour écrire correctement Schwinger--Dyson par rapport à Haar, et ne pas feindre que la
mesure de référence est plate, soit $\lambda_m$ le produit de Haar, de Lebesgue et du
dénombrement uniforme de la carte relevée, et soit
$R_m:=d\widehat\mu_m^0/d\lambda_m>0$. L’action totale régulée est
\begin{equation}\label{eq:pdf2-ym-total-regulated-action}
 \mathscr S_{m,g}:=S_{m,g}-\log R_m,
 \qquad
 d\widehat\mu_{m,g}^{\rm YM}
 =\mathscr Z_{m,g}^{-1}e^{-\mathscr S_{m,g}}d\lambda_m.
\end{equation}
Le terme $-\log R_m$ contient les noyaux de chaleur et les cartes de pont de la
structure de jauge de départ ; il est indépendant de $g$. Le terme de Yang--Mills qui
change le couplage et dont la limite classique est calculée ci-dessous est exactement
$S_{m,g}=\mathcal P_m^{F^2}/(4g^2)$.
La cohérence de la seconde égalité et le caractère standard des espaces finis
permettent d’appliquer Kolmogorov et définissent l’unique mesure cylindrique
$\widehat\mu_{\infty,g}^{\rm YM}$ dont les marginales sont
$\widehat\mu_{m,g}^{\rm YM}$.
Aucune valeur propre n’a été utilisée ici : $g$, l’action et
la mesure ont d’abord été fixés.

Pour définir la dynamique sur cette mesure sans perdre le produit, la jauge ni les
identités de surface, on construit le transport, on ne le postule pas. Dans chaque
facteur non atomique $X_{m,b}$, on fixe une carte borélienne
$\rho_{m,b}:X_{m,b}/\mathcal G_{m,b}\to(0,1)$ et on désintègre d’abord sur le
quotient de jauge, puis sur l’orbite de Haar. Soient $F_{m,b}^0$ et
$F_{m,b}^g$ les fonctions de répartition continues de $\rho_{m,b}$ pour la
mesure de référence et la mesure repondérée. Le transport triangulaire est
\begin{equation}\label{eq:pdf2-ym-gibbs-transport-factor}
 t_{m,b,g}:=(F_{m,b}^g)^{-1}\circ F_{m,b}^0,
 \qquad
 T_{m,b,g}(y,h):=(t_{m,b,g}(y),h),
 \qquad
 (T_{m,b,g})_*\lambda_{m,b}^0=\lambda_{m,b}^g.
\end{equation}
Dans un facteur purement atomique sur lequel $p_{m,b}=0$, on prend $T_{m,b,g}=I$. En présence d'atomes et d'une coordonnée continue, la carte est l'ordre lexicographique atome--intervalle et la même formule s'applique à l'espace standard non atomique total. Les transports des $9\cdot81$ couleurs se composent dans l'ordre fixé par \eqref{eq:pdf2-ym-owner-729-coloring} ; à l'intérieur d'une couleur, ils commutent car les colliers sont disjoints. On obtient
\begin{equation}\label{eq:pdf2-ym-gibbs-transport-global}
 T_{m,g}:=\prod_{r\in(\mathbb Z/9\mathbb Z)^3}
           \prod_{b\in\mathcal B_{m,r}^{F}}T_{m,b,g},
 \quad
 (T_{m,g})_*\widehat\mu_m^0=\widehat\mu_{m,g}^{\rm YM},
 \quad
 \widehat\pi_{m+1,m}T_{m+1,g}=T_{m,g}\widehat\pi_{m+1,m}.
\end{equation}
La dernière égalité se vérifie facteur par facteur en utilisant
\eqref{eq:pdf2-ym-action-cocycle} : le transport du facteur grossier est celui déjà fixé et
les nouveaux transports agissent uniquement sur l’innovation. Les cartes de quotient sont
choisies sur un écran représentant et transportées ; c’est pourquoi $T_{m,g}$ commute avec
la jauge et la réflexion et, pour $a\in E(4)$, vérifie
$T_{a\alpha,g}=U_\alpha(a)T_{\alpha,g}U_\alpha(a)^{-1}$.

La composition
\begin{equation}\label{eq:pdf2-ym-full-probability-isomorphism}
 \mathcal T_{m,g}:L^\infty(\widehat X_m,\widehat\mu_{m,g}^{\rm YM})
 \longrightarrow L^\infty(\widehat X_m,\widehat\mu_m^0),
 \qquad \mathcal T_{m,g}F:=F\circ T_{m,g}
\end{equation}
est, par construction, un isomorphisme d'espaces de probabilité et de $*$-algèbres, et se prolonge unitairement à $L^2$. En particulier,
\begin{equation}\label{eq:pdf2-ym-full-algebra-identities}
 \mathcal T_{m,g}(FG)=\mathcal T_{m,g}F\,\mathcal T_{m,g}G,
 \quad
 \mathcal T_{m,g}(F^*)=(\mathcal T_{m,g}F)^*,
 \quad
 \int F\,d\widehat\mu_{m,g}^{\rm YM}
 =\int\mathcal T_{m,g}F\,d\widehat\mu_m^0.
\end{equation}
Le domaine de chaque $T_{m,b,g}$ est l’espace même des holonomies qui satisfait
la loi de subdivision et la contrainte de Bianchi. Son image demeure donc
dans cet espace. En appliquant la multiplicativité précédente générateur par générateur, on
obtient les identités ponctuelles
\begin{equation}\label{eq:pdf2-ym-transport-surface-bianchi}
 \begin{aligned}
 \mathcal T_{m,g}(\operatorname{Hol}_p)
 &=\prod_{p'\subset p}^{\longrightarrow}
   \mathcal T_{m,g}(T_{p'}\operatorname{Hol}_{p'}T_{p'}^{-1}),\\
 \mathcal T_{m,g}(F_{{\rm cl},m})
 &=F_{{\rm cl},m}\circ T_{m,g},\\
 \prod_{p\subset\partial c}^{\longrightarrow}
 \mathcal T_{m,g}(T_p\operatorname{Hol}_pT_p^{-1})&=e,
 \qquad
 \operatorname{Alt}_{\lambda\mu\nu}
 D_\lambda(F_{{\rm cl},m}\circ T_{m,g})_{\mu\nu}=0.
 \end{aligned}
\end{equation}
La troisième ligne est l'annulation arête par arête de la frontière orientée de la frontière de $c$ ; la quatrième est sa polarisation infinitésimale. Ainsi, le transport préserve les produits, l'état, le clover et Bianchi, et pas seulement le premier chaos.

Enfin, on transporte les espérances, le générateur, le noyau, la projection et l’inclusion :
\begin{equation}\label{eq:pdf2-ym-g-dynamics}
 \begin{aligned}
 E_{m,\iota}^{g}&:=\mathcal T_{m,g}^{-1}E_{m,\iota}\mathcal T_{m,g},&
 \widehat H_{m,g}^{\rm YM}&:=\mathcal T_{m,g}^{-1}\widehat H_m^0\mathcal T_{m,g},\\
 \widehat K_{m,t}^{g}&:=\mathcal T_{m,g}^{-1}\widehat K_{m,t}^{0}\mathcal T_{m,g},&
 \mathbb P_{m,g}^{\rm phys}&:=\mathcal T_{m,g}^{-1}
     \mathbb P_m^{\rm phys}\mathcal T_{m,g},\\
 \widehat J_{m,g}&:=\mathcal T_{m+1,g}^{-1}
     \widehat J_m\mathcal T_{m,g},&
 D_{m,g}^{\rm YM}&:=\left.\widehat H_{m,g}^{\rm YM}
     \right|_{\operatorname{Ran}\mathbb P_{m,g}^{\rm phys}}.
 \end{aligned}
\end{equation}
La première ligne n’est pas seulement une conjugaison abstraite. En utilisant
\eqref{eq:pdf2-ym-action-factor-measures}, son action sur un cylindre est
l’actualisation par bain thermique de l’action :
\begin{equation}\label{eq:pdf2-ym-gibbs-conditional-expectation}
 (E_{m,b}^{g}F)(x_{\ne b})
 =\frac{\displaystyle
   \int_{X_{m,b}}F(x_{\ne b},u)e^{-p_{m,b}(u)/(4g^2)}
       d\lambda_{m,b}^0(u)}{\displaystyle
   \int_{X_{m,b}}e^{-p_{m,b}(u)/(4g^2)}d\lambda_{m,b}^0(u)},
 \qquad
 \widehat H_{m,g}^{\rm YM}
 =3\kappa_{\rm av}\sum_b(I-E_{m,b}^{g}).
\end{equation}
Ainsi, tant la mesure que chaque transition locale contiennent $g$ et l’action
$p_{m,b}/(4g^2)$ de manière explicite.
Elles sont markoviennes parce que $\mathcal T_{m,g}\mathbf1=\mathbf1$, locales parce que chaque
$T_{m,b,g}$ a le même collier que $p_{m,b}$, et dépendent de $g$ par
\eqref{eq:pdf2-ym-gibbs-transport-factor}. L’équivalence est spectrale et ne
sélectionne pas la mesure à partir du spectre :
Pour fermer également les équations dynamiques, soit $X_{m,b}^a$ le champ
invariant à gauche dans le facteur de lien $b$ associé à
$e_a\in\mathfrak g$ ; dans une carte continue d’innovation, on utilise la dérivée
directionnelle ordinaire, et les coordonnées finies sont traitées par différence exacte.
Pour tout cylindre lisse $F$, le domaine commun est
$\mathscr C_m^1:=\bigcap_{b,a}\operatorname{Dom}X_{m,b}^a$, en prenant un support
compact dans les cartes ouvertes d’innovation ; dans les facteurs de groupe, il n’y a pas de
frontière. L’invariance de Haar et la règle du produit calculent
\begin{equation}\label{eq:pdf2-ym-schwinger-dyson-finite}
 \int X_{m,b}^aF\,d\widehat\mu_{m,g}^{\rm YM}
 =\int F\,X_{m,b}^a\mathscr S_{m,g}\,d\widehat\mu_{m,g}^{\rm YM},
 \qquad F\in\mathscr C_m^1.
\end{equation}
Dans une coordonnée finie, pour $\Delta_\varepsilon F(q):=F(q+\varepsilon)-F(q)$, le changement de variable $q\mapsto q+\varepsilon$ donne la version exacte, sans dérivée formelle,
\begin{equation}\label{eq:pdf2-ym-schwinger-dyson-discrete}
 \int\Delta_\varepsilon F(q)\,d\widehat\mu_{m,g}^{\rm YM}
 =\int F(q)
  \left(e^{\mathscr S_{m,g}(q)-\mathscr S_{m,g}(q-\varepsilon)}-1\right)
  d\widehat\mu_{m,g}^{\rm YM}(q).
\end{equation}
Si $\delta_\varepsilon$ est la somme orientée de ces champs sur les arêtes qui
sont incidentes à un sommet, l’invariance de jauge de la mesure de référence et de
$S_{m,g}$ donne $\delta_\varepsilon\mathscr S_{m,g}=0$ et, par conséquent, l’identité de Ward
\begin{equation}\label{eq:pdf2-ym-ward-finite}
 \int\delta_\varepsilon F\,d\widehat\mu_{m,g}^{\rm YM}=0.
\end{equation}
Pour le facteur de jauge fini $\mathbb F_3^4$, l'invariance $\mathscr S_{m,g}(q+\varepsilon)=\mathscr S_{m,g}(q)$ réduit \eqref{eq:pdf2-ym-schwinger-dyson-discrete} à la même identité de Ward avec $\delta_\varepsilon$ remplacée par $\Delta_\varepsilon$. Les deux expressions sont compatibles avec $\widehat J_{m,g}$ car \eqref{eq:pdf2-ym-action-cocycle} sépare les dérivées grossières de celles du détail. Tout cylindre de la limite réside à un certain niveau ; ainsi, \eqref{eq:pdf2-ym-schwinger-dyson-finite} et \eqref{eq:pdf2-ym-ward-finite} passent littéralement à l'état cofinal. Ni l'équation de Schwinger--Dyson ni celle de Ward ne contiennent l'écart spectral comme prémisse.

\begin{equation}\label{eq:pdf2-ym-g-gap-before-limit}
 D_{m,g}^{\rm YM}\succeq3\kappa_{\rm av}
 (I-P_{\Omega_{m,g}}),
 \qquad
 W_{c,g}:=\mathcal T_{m,g}^{-1}W_c,
 \qquad
 D_{m,g}^{\rm YM}W_{c,g}=3\kappa_{\rm av}W_{c,g}.
\end{equation}
Dans la suite, on fixe un $g\in(0,\infty)$ et on supprime cet indice ; les
mesures, espérances, inclusions et générateurs sans exposant sont les objets
de \eqref{eq:pdf2-ym-g-dynamics}. Les holonomies géométriques demeurent des
fonctions coordonnées canoniques sur l’espace des configurations et, de ce fait, leur
loi change bien avec $g$ ; $\mathcal T_{m,g}$ les transporte en fonctions composées de la
carte de référence. Seul le témoin spectral est transporté explicitement dans
\eqref{eq:pdf2-ym-g-gap-before-limit}. Cette convention n’efface pas $g$ : sa résidence explicite est
\eqref{eq:pdf2-ym-gibbs-measure}--\eqref{eq:pdf2-ym-g-dynamics}.

\section{Vide simple et témoin qui atteint l'écart spectral}

L'intersection de toutes les espérances de l'état complet est la constante :
\begin{equation}\label{eq:pdf2-ym-owner-vacuum}
 \bigcap_{\iota\in\mathcal I_m}\operatorname{Ran}E_{m,\iota}
 =\mathbb C\Omega_m.
\end{equation}
Cette identité n'est pas introduite comme hypothèse. Dans la carte produit \eqref{eq:pdf2-ym-owner-product-factorization}, chaque $E_{m,\iota}$ intègre une coordonnée de l'indice exhaustif \eqref{eq:pdf2-ym-owner-exhaustive-index} et laisse les autres fixes. Si $u$ appartient à l'intersection, ses espérances conditionnelles par rapport à tout cylindre fini sont constantes. La convergence des martingales de ces cylindres vers la $\sigma$-algèbre complète implique que $u$ est constante presque partout. L'inclusion inverse est immédiate car toute espérance conserve l'unité. Cela prouve \eqref{eq:pdf2-ym-owner-vacuum} et, après la projection de jauge, laisse un vide physique unidimensionnel.
La décomposition de Hoeffding de \eqref{eq:pdf2-ym-owner-base-generator}--\eqref{eq:pdf2-ym-owner-refined-generator} donne
\begin{equation}\label{eq:pdf2-ym-owner-hoeffding}
 \widehat{\mathcal E}_m(u)
 =3\kappa_{\rm av}\sum_{S\Subset\mathcal I_m}|S|\,\|u_S\|^2,
 \qquad
 D_m^{\rm YM}\succeq
 3\kappa_{\rm av}(I-P_{\Omega_m}).
\end{equation}
La fonction
\begin{equation}\label{eq:pdf2-ym-owner-wc}
 W_c(q):=\mathbf1_{\{q_1+\cdots+q_6=0\}}-\frac13
\end{equation}
est invariant sous $G^{V(\Lambda_m)}$ car il ne dépend ni des liens ni des repères. Pour
l’action d’incidence, chaque $x_i$ apparaît trois fois, donc
\begin{equation}\label{eq:pdf2-ym-owner-wc-invariance}
 \sum_{i<j}(B^+x)_{ij}=3\sum_i x_i=0
 \quad\text{dans }\mathbb F_3.
\end{equation}
Ainsi $\mathbb P_m^{\rm phys}W_c=W_c$. Sous $u_6$, la somme des six trits est
uniforme ; par calcul direct,
\begin{equation}\label{eq:pdf2-ym-owner-wc-moments}
 \mathbb EW_c=0,\qquad
 \|W_c\|^2=\frac29,\qquad
 \mathbb E(W_c^3)=\frac2{27}.
\end{equation}
Seul le projecteur $I-E_{m,q_c}$ agit non trivialement sur $W_c$. Par conséquent
\begin{equation}\label{eq:pdf2-ym-owner-wc-eigenvalue}
 D_m^{\rm YM}W_c=3\kappa_{\rm av}W_c.
\end{equation}
La borne inférieure de \eqref{eq:pdf2-ym-owner-hoeffding} n’est pas seulement une borne : le
complément physique contient un vecteur non nul qui atteint le seuil.

\section{Raffinement nonadique et persistance dans la limite}

Les inclusions physiques satisfont
\begin{equation}\label{eq:pdf2-ym-owner-intertwining}
 D_{m+1}^{\rm YM}\widehat J_m
 =\widehat J_mD_m^{\rm YM},\qquad
 \bigl(e^{-a_{m+1}D_{m+1}^{\rm YM}}\bigr)^9\widehat J_m
 =\widehat J_me^{-a_mD_m^{\rm YM}},\qquad a_{m+1}=a_m/9.
\end{equation}
Dans la limite inductive
\begin{equation}\label{eq:pdf2-ym-owner-inductive-limit}
 \widehat{\mathscr H}_\infty^{\rm phys}
 =\varinjlim(\widehat{\mathscr H}_m^{\rm phys},\widehat J_m)
\end{equation}
est défini sur l’union algébrique
\begin{equation}\label{eq:pdf2-ym-owner-limit-form}
 \mathcal E_\infty[\widehat J_{m,\infty}u]
 :=\langle u,D_m^{\rm YM}u\rangle.
\end{equation}
Son générateur est noté
$D_{\infty,g}^{\rm YM}$ ; conformément à la convention postérieure à
\eqref{eq:pdf2-ym-g-gap-before-limit}, on écrit aussi
$D_\infty^{\rm YM}:=D_{\infty,g}^{\rm YM}$.
Le premier entrelacement de \eqref{eq:pdf2-ym-owner-intertwining} rend la
définition indépendante du représentant. Les troncatures de Hoeffding sont
un cœur de forme et donnent une suite de récupération ; Fatou sur la somme non négative
donne la semi-continuité inférieure. La fermeture possède donc
\begin{equation}\label{eq:pdf2-ym-owner-limit-gap}
 \overline{\mathcal E}_\infty[u]
 \ge3\kappa_{\rm av}\|(I-P_{\Omega_\infty})u\|^2.
\end{equation}
L’inclusion conserve les coordonnées ancestrales, de sorte que
$\widehat J_{m,\infty}W_c$ reste non nul et vérifie
\begin{equation}\label{eq:pdf2-ym-owner-limit-witness}
 D_\infty^{\rm YM}\widehat J_{m,\infty}W_c
 =3\kappa_{\rm av}\widehat J_{m,\infty}W_c.
\end{equation}
En conséquence,
\begin{equation}\label{eq:pdf2-ym-owner-exact-gap}
 \inf\bigl(\operatorname{Spec}D_\infty^{\rm YM}\setminus\{0\}\bigr)
 =3\kappa_{\rm av}>0,
 \qquad \ker D_\infty^{\rm YM}=\mathbb C\Omega_\infty.
\end{equation}

\section{Quatre reconstructions OS de la même dynamique}

Le noyau complet de \eqref{eq:pdf2-ym-owner-full-kernel} définit une loi réversible sur $\widehat X_m$. Pour la reconstruction physique, on utilise sa descente bien typée \eqref{eq:pdf2-ym-owner-physical-kernel} : sur le quotient $Y_m$, on définit
\begin{equation}\label{eq:pdf2-ym-owner-eight-face-law}
 \Omega_{m,{\rm phys}}^{(8)}(d\bar y_1\cdots d\bar y_8)
 =\int\nu_m(d\bar z)
  \prod_{r=1}^8K_{m,a_m/2}^{\rm phys}(\bar z,d\bar y_r).
\end{equation}
Pour chaque direction $\nu=1,\ldots,4$, la réflexion échange les quatre faces
positives avec les négatives. En conditionnant par rapport à l’état central $z$ et en utilisant
la réversibilité,
\begin{equation}\label{eq:pdf2-ym-owner-four-os}
 \int\overline{F(\Theta_{\nu,m}y)}F(y)\,
 d\Omega_{m,{\rm phys}}^{(8)}(y)
 =\|\mathcal B_{m,\nu}F\|_2^2\ge0.
\end{equation}
La marginale de la paire opposée normale à $\nu$ est
\begin{equation}\label{eq:pdf2-ym-owner-os-marginal}
 \nu_m(d\bar x)K_{m,a_m}^{\rm phys}(\bar x,d\bar y).
\end{equation}
Pour identifier l'hamiltonien d'une direction, on prend le sous-espace OS des observables d'une face positive simple, $F(y)=F_\nu(\bar y_{\nu,+})$ ; les observables de plusieurs faces conservent la positivité précédente, mais ne sont pas confondus avec cet espace de transfert. Dans le quotient de face simple, l'application
\begin{equation}\label{eq:pdf2-ym-owner-os-unitary}
 \mathcal V_{m,\nu}[F]
 :=e^{-a_mD_m^{\rm YM}/2}\mathcal W_mF
\end{equation}
est isométrique pour la norme OS et a une image dense. Son extension identifie le
quotient OS à $\widehat{\mathscr H}_m^{\rm phys}$. Posons
$\overline D_m:=\mathcal W_m^{-1}D_m^{\rm YM}\mathcal W_m$ et
$\overline J_m:=\mathcal W_{m+1}^{-1}\widehat J_m\mathcal W_m$.
L’entrelacement et $a_m=9a_{m+1}$ donnent
\begin{equation}\label{eq:pdf2-ym-owner-os-range-inclusion}
 \overline J_me^{-a_m\overline D_m/2}
 =e^{-a_m\overline D_{m+1}/2}\overline J_m
 =\bigl(e^{-a_{m+1}\overline D_{m+1}/2}\bigr)^9\overline J_m.
\end{equation}
Le dernier membre appartient à l’image de
$e^{-a_{m+1}\overline D_{m+1}/2}$ ; c’est pourquoi l’inverse de $\mathcal V_{m+1,\nu}$
qui apparaît ci-après ne s’applique que sur son image et est bien défini. Les
connecteurs
\begin{equation}\label{eq:pdf2-ym-owner-os-connectors}
 \mathcal J_{m,\nu}^{\rm OS}
 :=\mathcal V_{m+1,\nu}^{-1}\widehat J_m\mathcal V_{m,\nu}
\end{equation}
sont des isométries et entrelacent les quatre hamiltoniens reconstruits :
\begin{equation}\label{eq:pdf2-ym-owner-common-os-hamiltonian}
 H_{{\rm OS},m+1}^{(\nu)}\mathcal J_{m,\nu}^{\rm OS}
 =\mathcal J_{m,\nu}^{\rm OS}H_{{\rm OS},m}^{(\nu)},
 \qquad
 \mathcal V_{\infty,\nu}H_{{\rm OS},\infty}^{(\nu)}
 \mathcal V_{\infty,\nu}^{-1}=D_\infty^{\rm YM}.
\end{equation}
Les quatre réflexions ne fabriquent pas quatre mesures : ce sont quatre publications du même
état, du même noyau et du même générateur.

\section{Localité, action euclidienne et publication de courbure}

La catégorie dirigée des écrans comprend les raffinements, les superpositions et les transports
euclidiens. Si $a\in E(4)$ et $\alpha$ est un écran, le réétiquetage complet définit
\begin{equation}\label{eq:pdf2-ym-owner-e4-finite-transport}
 U_\alpha(a):\widehat{\mathscr H}_\alpha^{\rm phys}
 \longrightarrow\widehat{\mathscr H}_{a\alpha}^{\rm phys},
 \qquad
 U_\beta(a)\widehat J_{\alpha\beta}
 =\widehat J_{a\alpha,a\beta}U_\alpha(a).
\end{equation}
Le transport préserve la mesure, le générateur, les idéaux nuls et les supports. L'identité de naturalité permet de définir, sans choisir de représentant, une application unitaire sur la colimite algébrique :
\begin{equation}\label{eq:pdf2-ym-owner-e4-colimit-action}
 U(a)\widehat J_{\alpha,\infty}F
 :=\widehat J_{a\alpha,\infty}U_\alpha(a)F.
\end{equation}
Les réétiquetages vérifient $U_{a\alpha}(b)U_\alpha(a)=U_\alpha(ba)$ et préservent
la norme ; c’est pourquoi \eqref{eq:pdf2-ym-owner-e4-colimit-action} se prolonge en une
représentation unitaire de $E(4)$.

\subsection{Le canal d'incidence est une courbure locale, non un facteur auxiliaire}

Tout groupe compact, connexe et simple possède un tore maximal non trivial. On fixe $t_3\in G$ d'ordre exactement trois. Pour un collier $c$, soit
\begin{equation}\label{eq:pdf2-ym-owner-incidence-flux}
 \omega_c(q):=\sum_{i<j}q_{ij}\in\mathbb F_3.
\end{equation}
L’identité \eqref{eq:pdf2-ym-incidence-conservation} prouve que $\omega_c$ est
invariant sous l’action finie. Si $h_c$ est le repère au sommet de base du collier,
on définit l’holonomie géométrique et, pour le témoin spectral, son image transportée
dans la loi au couplage $g$ :
\begin{equation}\label{eq:pdf2-ym-owner-incidence-holonomy}
 \operatorname{Hol}_{m,c}^{\rm inc,geom}:=h_ct_3^{\,\omega_c(q)}h_c^{-1}\in G,
 \qquad
 \operatorname{Hol}_{m,c,g}^{\rm inc,spec}
 :=\mathcal T_{m,g}^{-1}\operatorname{Hol}_{m,c}^{\rm inc,geom},
 \qquad
 \operatorname{Hol}_{m,c}^{\rm inc}:=\operatorname{Hol}_{m,c,g}^{\rm inc,spec}.
\end{equation}
Sous $G^{V_m}$, elle se transforme par conjugaison et sous
$\prod_c\mathbb F_3^4$, elle demeure fixe. Comme $t_3$ est d’ordre trois, le calcul
fonctionnel de cette variable à trois valeurs donne l’identité littérale
\begin{equation}\label{eq:pdf2-ym-owner-wc-curvature-functional}
 W_{c,g}:=\mathcal T_{m,g}^{-1}W_c^0
 =\mathbf1_{\{e\}}\!\left(\operatorname{Hol}_{m,c}^{\rm inc}\right)-\frac13.
\end{equation}
On écrit $W_c:=W_{c,g}$ dans le reste du chapitre. Comme $T_{m,b,g}$ agit à l'intérieur du quotient du même collier, l'holonomie spectrale est une fonction de son algèbre locale d'holonomies géométriques ; elle n'ajoute pas de variable et ne change pas le support. Les holonomies géométriques, utilisées dans $S_{m,g}$ et dans le clover, maintiennent leur distribution dépendante de $g$.
Le projecteur du membre de droite est une fonction de classe locale ; par conséquent, $W_c$ n'est pas une coordonnée auxiliaire découplée, mais un observable de l'holonomie d'incidence de la même restriction de jauge.

Pour un ouvert borné $O\subset\mathbb R^4$, soit
$\mathfrak A_{\rm YM}(O)$ l’algèbre engendrée par les fonctions de classe bornées de
\eqref{eq:pdf2-ym-owner-incidence-holonomy}, les holonomies de plaquette et les lecteurs
clover à support dans $O$. On définit le secteur de courbure physique par
\begin{equation}\label{eq:pdf2-ym-owner-curvature-cyclic-sector}
 \mathscr H_{\rm YM}^{\rm curv}
 :=\overline{\bigcup_{O\Subset\mathbb R^4}
   \mathfrak A_{\rm YM}(O)\Omega_\infty}
 \subseteq\widehat{\mathscr H}_\infty^{\rm phys}.
\end{equation}
Par \eqref{eq:pdf2-ym-owner-wc-curvature-functional}, chaque
$\widehat J_{m,\infty}W_c$ appartient matériellement à ce secteur.
Les espérances de coordonnée envoient les cylindres de courbure sur des cylindres de
courbure ; le semi-groupe laisse donc invariant
$\mathscr H_{\rm YM}^{\rm curv}$. Étant auto-adjoint, le secteur est réducteur. La
borne de \eqref{eq:pdf2-ym-owner-limit-gap} s’y restreint et le vecteur précédent
atteint l’égalité :
\begin{equation}\label{eq:pdf2-ym-owner-curvature-sector-gap}
 D_{\rm curv,g}^{\rm YM}
 :=\left.D_{\infty,g}^{\rm YM}\right|_{\mathscr H_{\rm YM}^{\rm curv}},
 \qquad D_{\rm curv}^{\rm YM}:=D_{\rm curv,g}^{\rm YM},
 \qquad
 \inf\bigl(\operatorname{Spec}D_{\rm curv,g}^{\rm YM}\setminus\{0\}\bigr)
 =3\kappa_{\rm av}.
\end{equation}

\subsection{Système cofinal de publication et produits croisés}

L’inclusion $\widehat J_{\alpha\beta}$ conserve une coordonnée ancestrale ; elle ne doit pas
être confondue avec la moyenne de bloc qui publie un champ continu. Pour écrire
cette seconde application sans altérer la généalogie, soit $\mathfrak P$ l’ensemble dirigé
des écrans cellulaires bornés, fermé sous superposition, raffinement nonadique et
transport euclidien. Pour $\alpha\in\mathfrak P$, posons
\begin{equation}\label{eq:pdf2-ym-owner-cell-step-space}
 V_\alpha:=\operatorname{span}\left\{
 e_{\alpha,c}:=|c|^{-1/2}\mathbf1_c:
 c\in\mathcal C_\alpha^{\rm cell}\right\}
 \subset L^2(\mathbb R^4),
 \qquad P_\alpha:L^2(\mathbb R^4)\to V_\alpha.
\end{equation}
Les superpositions font que $P_\alpha h\to h$ pour tout $h\in L^2$ le long de
$\mathfrak P$. Du côté de la jauge, on normalise le témoin par
\begin{equation}\label{eq:pdf2-ym-owner-normalized-cell-witness}
 Z_{\alpha,c}:=\sqrt{\frac92}\,W_{\alpha,c},
 \qquad
 \langle Z_{\alpha,c},Z_{\alpha,d}\rangle=\delta_{cd},
 \qquad
 D_\alpha^{\rm YM}Z_{\alpha,c}=3\kappa_{\rm av}Z_{\alpha,c}.
\end{equation}
La seconde égalité utilise exactement les moments de \eqref{eq:pdf2-ym-owner-wc-moments} et la factorisation de coordonnées distinctes. Soit $\mathcal K_\alpha^{\rm cell}$ l'espace engendré par ces vecteurs. Si $\beta\succeq\alpha$, l'opérateur d'entrelacement de publication, distinct de $\widehat J_{\alpha\beta}$, est défini générateur par générateur par
\begin{equation}\label{eq:pdf2-ym-owner-block-intertwiner}
 I_{\alpha\beta}^{\rm curv}Z_{\alpha,c}
 :=\sum_{\substack{d\in\mathcal C_\beta^{\rm cell}\\d\subset c}}
       \sqrt{\frac{|d|}{|c|}}\,Z_{\beta,d}.
\end{equation}
Comme les sous-cellules sont disjointes et ont pour somme $|c|$, on obtient, sans passage à la limite,
\begin{equation}\label{eq:pdf2-ym-owner-block-intertwiner-identities}
 (I_{\alpha\beta}^{\rm curv})^*I_{\alpha\beta}^{\rm curv}=I,
 \quad
 I_{\beta\gamma}^{\rm curv}I_{\alpha\beta}^{\rm curv}
 =I_{\alpha\gamma}^{\rm curv},
 \quad
 D_\beta^{\rm YM}I_{\alpha\beta}^{\rm curv}
 =I_{\alpha\beta}^{\rm curv}D_\alpha^{\rm YM}.
\end{equation}
Le réétiquetage cellulaire conserve les volumes et les descendants ; pour $a\in E(4)$, on a en outre
\begin{equation}\label{eq:pdf2-ym-owner-block-e4-naturality}
 U_\beta(a)I_{\alpha\beta}^{\rm curv}
 =I_{a\alpha,a\beta}^{\rm curv}U_\alpha(a).
\end{equation}

Pour $h\in C_c^\infty(\mathbb R^4)$, définissons maintenant le lecteur de bloc
\begin{equation}\label{eq:pdf2-ym-owner-smeared-witness}
 \Phi_\alpha(h):=
 \sum_{c\in\mathcal C_\alpha^{\rm cell}}
 \langle h,e_{\alpha,c}\rangle_{L^2}Z_{\alpha,c}.
\end{equation}
Sur un maillage uniforme, cette formule est
$\sqrt{9/2}\,a_\alpha^2\sum_c(P_\alpha h)(x_c)W_{\alpha,c}$ ; la moyenne cellulaire,
non une évaluation ponctuelle, fixe le produit croisé. En effet, si $\gamma$ raffine
deux écrans $\alpha$ et $\beta$, les définitions précédentes donnent l’identité
calculée
\begin{equation}\label{eq:pdf2-ym-owner-smeared-cross-gram}
 \left\langle
  I_{\alpha\gamma}^{\rm curv}\Phi_\alpha(h),
  I_{\beta\gamma}^{\rm curv}\Phi_\beta(k)
 \right\rangle
 =\langle P_\alpha h,P_\beta k\rangle_{L^2}.
\end{equation}
Par polarisation, avec $k=h$,
\begin{equation}\label{eq:pdf2-ym-owner-smeared-cauchy-derived}
 \left\|
  I_{\alpha\gamma}^{\rm curv}\Phi_\alpha(h)
  -I_{\beta\gamma}^{\rm curv}\Phi_\beta(h)
 \right\|_2^2
 =\|P_\alpha h-P_\beta h\|_{L^2}^2\longrightarrow0.
\end{equation}
Ainsi, la propriété de Cauchy est une conséquence de la convergence forte des
projections, et non une identité supplémentaire postulée. La colimite de
\eqref{eq:pdf2-ym-owner-block-intertwiner} contient donc
\begin{equation}\label{eq:pdf2-ym-owner-smeared-gram}
 \Phi(h):=\lim_\alpha I_{\alpha,\infty}^{\rm curv}\Phi_\alpha(h),
 \qquad
 \langle\Phi(h),\Phi(k)\rangle=\langle h,k\rangle_{L^2},
 \qquad
 D_{\rm pub}^{\rm YM}\Phi(h)=3\kappa_{\rm av}\Phi(h).
\end{equation}

Cette publication n'ajoute pas de degrés de liberté. À chaque étape, la décomposition $V_\beta=V_\alpha\oplus(V_\beta\ominus V_\alpha)$ possède une base de Haar nonadique. L'ordre exhaustif des nouvelles incidences de \eqref{eq:pdf2-ym-owner-exhaustive-index} apparie cette base, étiquette par étiquette, avec des coordonnées $Z_{\beta,d}$ du facteur d'innovation, en conservant le support de la fonction de Haar. Les applications unitaires triangulaires $\widehat\Gamma_\alpha$ recollent ces appariements et produisent une isométrie basée
\begin{equation}\label{eq:pdf2-ym-owner-publication-into-physical}
 \mathcal U_{\rm curv}:
 \varinjlim(\mathcal K_\alpha^{\rm cell},I_{\alpha\beta}^{\rm curv})
 \longrightarrow\mathscr H_{\rm YM}^{\rm curv},
 \quad
 D_{\rm curv}^{\rm YM}\mathcal U_{\rm curv}
 =\mathcal U_{\rm curv}D_{\rm pub}^{\rm YM},
 \quad U(a)\mathcal U_{\rm curv}=\mathcal U_{\rm curv}U_{L^2}(a).
\end{equation}
Désormais, $\Phi(h)$ désigne également son image par $\mathcal U_{\rm curv}$.
La compatibilité de \eqref{eq:pdf2-ym-gibbs-transport-global} définit
$\mathcal T_{\infty,g}$ à la limite. La publication physique au couplage fixé
n’est pas une seconde isométrie choisie uniquement sur le premier chaos, mais la restriction
\begin{equation}\label{eq:pdf2-ym-full-publication-restriction}
 \mathcal U_{{\rm curv},g}
 :=\mathcal T_{\infty,g}^{-1}\mathcal U_{{\rm curv},0},
 \qquad
 \widetilde{\mathcal U}_g
 :=\mathcal T_{\infty,g}^{-1}:
 L^\infty(\widehat X_\infty,\widehat\mu_\infty^0)
 \longrightarrow
 L^\infty(\widehat X_\infty,\widehat\mu_{\infty,g}^{\rm YM}).
\end{equation}
Sur la $*$-algèbre polynomiale engendrée par les holonomies, les plaquettes, clover et
$\Phi(h)$, on vérifie, générateur par générateur puis par linéarité,
\begin{equation}\label{eq:pdf2-ym-full-publication-products}
 \widetilde{\mathcal U}_g(AB)
 =\widetilde{\mathcal U}_g(A)\widetilde{\mathcal U}_g(B),
 \quad
 \widetilde{\mathcal U}_g(A^*)=\widetilde{\mathcal U}_g(A)^*,
 \quad
 \omega_g(\widetilde{\mathcal U}_gA)=\omega_0(A).
\end{equation}
Les égalités de surface, de clover et de Bianchi sont préservées par
\eqref{eq:pdf2-ym-transport-surface-bianchi} ; par densité monotone, l’extension
vaut dans toute l’algèbre locale bornée. Ainsi,
\eqref{eq:pdf2-ym-owner-smeared-cross-gram} est la restriction d’un isomorphisme
probabiliste et algébrique complet, non un substitut de celui-ci.

\subsection{Noyau dense, continuité forte et réseau local}

Soit $M_{\Phi(h)}$ l'opérateur de multiplication affilié à l'algèbre locale et
\begin{equation}\label{eq:pdf2-ym-owner-curvature-weyl}
 \mathsf W(h,t):=\exp\!\bigl(itM_{\Phi(h)}\bigr),
 \qquad h\in C_c^\infty(\mathbb R^4),\quad t\in\mathbb R.
\end{equation}
On complète $\mathfrak A_{\rm YM}(O)$ par ces unitaires pour
$\operatorname{supp}h\subset O$ et par les fonctions de classe bornées des
holonomies de plaquette. Notons $\mathfrak A_{\rm YM}(O)^{\rm sm}$ la
$*$-algèbre ainsi engendrée avec des fonctions tests lisses. Le noyau
\begin{equation}\label{eq:pdf2-ym-owner-smooth-core}
 \mathscr D_{\rm sm}:=\operatorname{span}\left\{
 A_1\cdots A_n\Omega_\infty:
 A_j\in\mathfrak A_{\rm YM}(O_j)^{\rm sm},\quad
 O_j\Subset\mathbb R^4\right\}
\end{equation}
est dense par la définition cyclique de \eqref{eq:pdf2-ym-owner-curvature-cyclic-sector}. Les équations \eqref{eq:pdf2-ym-owner-block-e4-naturality} et \eqref{eq:pdf2-ym-owner-publication-into-physical} donnent, pour $a\in E(4)$,
\begin{equation}\label{eq:pdf2-ym-owner-field-e4-covariance}
 U(a)\Phi(h)=\Phi(a\!\cdot h),
 \qquad
 \|(U(a)-I)\Phi(h)\|=\|a\!\cdot h-h\|_{L^2}.
\end{equation}
En outre, $|e^{ix}-e^{iy}|\le|x-y|$ et
\eqref{eq:pdf2-ym-owner-smeared-gram} impliquent
\begin{equation}\label{eq:pdf2-ym-owner-weyl-continuity}
 \|\bigl(\mathsf W(a\!\cdot h,t)-\mathsf W(h,t)\bigr)\Omega_\infty\|
 \le |t|\,\|a\!\cdot h-h\|_{L^2}.
\end{equation}
La continuité se vérifie également pour les générateurs d'holonomie du cœur ; \eqref{eq:pdf2-ym-owner-weyl-continuity} à elle seule ne les inclut pas. Pour un chemin polygonal orienté $\gamma$, une plaquette $p$ et $f\in C^1(G)^{\rm class}$, on pose
\begin{equation}\label{eq:pdf2-ym-owner-holonomy-core-generators}
 A_{\gamma,f}:=f(\operatorname{Hol}_\gamma),
 \qquad A_{p,f}:=f(\operatorname{Hol}_p).
\end{equation}
Dans une superposition de $\gamma$ et de $a\gamma$, avec $a\in E(4)$, la loi de pont
\eqref{eq:pdf2-ym-owner-heat-bridge} annule les accroissements communs. Pour chaque
accroissement restant $z$, on utilise
$|f(xz)-f(x)|\le\|\nabla f\|_\infty d_G(z,e)$ et la borne du second moment du
noyau de chaleur
\(
 \int_Gd_G(z,e)^2k_\tau(z)\,dz\le C_G\tau
\).
Pour la loi repondérée, $0<e^{-p/(4g^2)}\le1$ et la continuité de $p$ sur le
facteur compact donnent
\begin{equation}\label{eq:pdf2-ym-owner-gibbs-heat-moment}
 \frac{\int_Gd_G(z,e)^2e^{-p(z)/(4g^2)}k_\tau(z)\,dz}
      {\int_Ge^{-p(z)/(4g^2)}k_\tau(z)\,dz}
 \le C_{G,g}\tau,
\end{equation}
uniformément sur les écrans d’un compact et pour $0<\tau\le\tau_0$.
En additionnant les temps des incréments de la différence symétrique, on obtient
l’estimation calculée
\begin{equation}\label{eq:pdf2-ym-owner-holonomy-e4-estimate}
 \begin{aligned}
 \|U(a)A_{\gamma,f}\Omega_\infty-A_{\gamma,f}\Omega_\infty\|_2^2
 &\le C_{G,g}\|\nabla f\|_\infty^2
     \tau(\gamma\triangle a\gamma),\\
 \|U(a)A_{p,f}\Omega_\infty-A_{p,f}\Omega_\infty\|_2^2
 &\le C_{G,g}\|\nabla f\|_\infty^2
     \tau(\partial p\triangle\partial(ap)).
 \end{aligned}
\end{equation}
Ici, la différence symétrique n'est pas mesurée comme deux courbes disjointes : les connecteurs de surface la remplissent par le ruban orienté minimal $\Sigma(\gamma,a\gamma)$ et
\begin{equation}\label{eq:pdf2-ym-owner-bridge-time-ribbon}
 \tau(\gamma\triangle a\gamma)
 :=\sum_{p\subset\Sigma(\gamma,a\gamma)}\tau_p,
 \qquad
 \tau(\gamma\triangle a\gamma)
 \le C_\gamma d_{E(4)}(a,e),
\end{equation}
avec la même définition pour les frontières de plaquette. L’inégalité découle de la
somme des temps additifs de \eqref{eq:pdf2-ym-owner-heat-bridge} sur une bande
de largeur $d_{E(4)}(a,e)$ et de longueur bornée par celle de $\gamma$.
La mesure du temps de pont du membre de droite tend ainsi vers zéro lorsque $a\to e$ ;
c’est exactement la continuité des connecteurs de route de l’image de jauge.
Chaque lecteur clover borné $\chi(F_{\rm cl})$, $\chi\in C_b^1$, est fonction d’une
somme finie de plaquettes transportées ; la même borne, multipliée par
$\|\chi'\|_\infty$ et par le nombre de termes, prouve également sa continuité.

Pour des fonctions de classe seulement bornées, on rend visible la régularisation qui
entre dans le cœur. Si $d b$ est la mesure de Haar à gauche sur $E(4)$ et
$\eta\in C_c^\infty(E(4))$, on définit l’intégrale de Bochner
\begin{equation}\label{eq:pdf2-ym-owner-orbit-smearing}
 A[\eta]:=\int_{E(4)}\eta(b)U(b)AU(b)^*\,db,
 \qquad A\in\{A_{\gamma,f},A_{p,f},\chi(F_{{\rm cl}})\},
 \quad \chi\in C_b^1.
\end{equation}
Avec $(L_a\eta)(b):=\eta(a^{-1}b)$, le changement de variable $b\mapsto ab$ donne
\begin{equation}\label{eq:pdf2-ym-owner-orbit-smearing-continuity}
 U(a)A[\eta]U(a)^*=A[L_a\eta],
 \qquad
 \|(A[L_a\eta]-A[\eta])\Omega_\infty\|
 \le\|A\|\,\|L_a\eta-\eta\|_{L^1(E(4))}\longrightarrow0.
\end{equation}
Il est désormais entendu que les générateurs d'holonomie, de plaquette et de clover de $\mathfrak A_{\rm YM}(O)^{\rm sm}$ sont ceux de \eqref{eq:pdf2-ym-owner-holonomy-core-generators} avec $f\in C^1$, ou leurs régularisations \eqref{eq:pdf2-ym-owner-orbit-smearing}, et que le support balayé reste compact à l'intérieur de $O$. Les approximations de l'identité dans $E(4)$, avec \eqref{eq:pdf2-ym-owner-holonomy-e4-estimate}, récupèrent les générateurs non régularisés dans $L^2$ ; le secteur cyclique n'est donc pas réduit.

Un télescopage de produits qui utilise
\eqref{eq:pdf2-ym-owner-weyl-continuity} et
\eqref{eq:pdf2-ym-owner-orbit-smearing-continuity} prouve la continuité de tous
les générateurs de $\mathscr D_{\rm sm}$ ; par
densité et unitarité,
\begin{equation}\label{eq:pdf2-ym-owner-strong-e4}
 \lim_{a\to e}\|(U(a)-I)\Psi\|=0
 \qquad(\Psi\in\mathscr H_{\rm YM}^{\rm curv}).
\end{equation}
Les supports disjoints utilisent des facteurs disjoints et commutent. On obtient ainsi le réseau
\begin{equation}\label{eq:pdf2-ym-owner-local-net}
\begin{aligned}
 O_1\subseteq O_2&\Rightarrow
 \mathfrak A_{\rm YM}(O_1)\subseteq\mathfrak A_{\rm YM}(O_2),\\
 [\mathfrak A_{\rm YM}(O_1),\mathfrak A_{\rm YM}(O_2)]&=0
 \quad (O_1\cap O_2=\varnothing),\\
 U(a)\mathfrak A_{\rm YM}(O)U(a)^*&=\mathfrak A_{\rm YM}(aO).
\end{aligned}
\end{equation}
L’identité de Gram pour $h\ne0$ et
\eqref{eq:pdf2-ym-owner-wc-curvature-functional} prouvent que le réseau n’est pas scalaire.

\subsection{Paquet OS continu et reconstruction}

Pour une direction $\nu$, soit $\mathfrak A_{+,\nu}$ l’adhérence de l’union des
algèbres précédentes à support compact dans $\{x_\nu>0\}$. L’état
statique, la loi des piles et l’opérateur de transfert ne sont pas définis
séparément. Sur un écran $\alpha$, soient
$\widehat\mu_{\alpha,g}^{\rm YM}$ la marginale de
\eqref{eq:pdf2-ym-gibbs-measure},
$\nu_{\alpha,g}:=(\pi_\alpha^{\rm phys})_*
\widehat\mu_{\alpha,g}^{\rm YM}$ sa loi sur le quotient et
$K_{\alpha,g}(t)=e^{-tD_{\alpha,g}^{\rm YM}}$ son noyau.
Pour des temps ordonnés $t_1\le\cdots\le t_n$, la marginale euclidienne de la
même mesure est
\begin{equation}\label{eq:pdf2-ym-owner-euclidean-time-marginal}
 \mathbb P_{\alpha,g}^{E,(\nu;t_1,\ldots,t_n)}
 :=(\operatorname{ev}_{t_1}^{(\nu)},\ldots,
    \operatorname{ev}_{t_n}^{(\nu)})_*
    \widehat\mu_{\alpha,g}^{\rm YM}.
\end{equation}
L'identité de désintégration qui fait partie de \eqref{eq:pdf2-g-gau-transfer-euclidean-dato} la calcule, sans la remplacer par une chaîne auxiliaire :
\begin{equation}\label{eq:pdf2-ym-owner-euclidean-markov-identification}
 d\mathbb P_{\alpha,g}^{E,(\nu;t_1,\ldots,t_n)}
 =\nu_{\alpha,g}(dy_1)
   \prod_{j=1}^{n-1}K_{\alpha,g}
      (t_{j+1}-t_j;y_j,dy_{j+1}).
\end{equation}
Si $M_A$ désigne la multiplication par un lecteur borné $A$, ses corrélations
sont donc
\begin{equation}\label{eq:pdf2-ym-owner-schwinger-transfer}
 \begin{aligned}
 S_{\alpha,n}^{(\nu,g)}(A_1,t_1;\ldots;A_n,t_n)
 &:=\int\prod_{j=1}^n(\tau_{t_je_\nu}A_j)(x)\,
       d\widehat\mu_{\alpha,g}^{\rm YM}(x)\\
 &=\langle\mathbf1,
 M_{A_1}K_{\alpha,g}(t_2-t_1)M_{A_2}\cdots
 K_{\alpha,g}(t_n-t_{n-1})M_{A_n}\mathbf1
 \rangle_{L^2(\nu_{\alpha,g})}.
 \end{aligned}
\end{equation}
La loi de la pile stationnaire à ces mêmes temps est cette marginale euclidienne, écrite au moyen de sa désintégration :
\begin{equation}\label{eq:pdf2-ym-owner-markov-stack-law}
 d\mathbb P_{\alpha,g}^{(\nu;t_1,\ldots,t_n)}
 :=d\mathbb P_{\alpha,g}^{E,(\nu;t_1,\ldots,t_n)}
 =\nu_{\alpha,g}(dy_1)
   \prod_{j=1}^{n-1}K_{\alpha,g}(t_{j+1}-t_j;y_j,dy_{j+1}).
\end{equation}
Intégrer successivement $y_n,y_{n-1},\ldots,y_1$ donne l’identité exacte
\begin{equation}\label{eq:pdf2-ym-owner-stack-transfer-identity}
 S_{\alpha,n}^{(\nu,g)}(A_1,t_1;\ldots;A_n,t_n)
 =\int\prod_{j=1}^nA_j(y_j)\,
   d\mathbb P_{\alpha,g}^{(\nu;t_1,\ldots,t_n)}(y).
\end{equation}
Si tous les temps coïncident, les noyaux sont l’identité et
\begin{equation}\label{eq:pdf2-ym-owner-static-is-stack}
 S_{\alpha,n}^{(\nu,g)}(A_1,t;\ldots;A_n,t)
 =\int A_1\cdots A_n\,d\nu_{\alpha,g}.
\end{equation}
Ainsi, l'état de Schwinger statique est la marginale à temps égal de la même mesure euclidienne, et sa désintégration est la loi de Markov ; ce ne sont pas deux fonctionnelles que l'on identifie ensuite par leur nom. La compatibilité de $\widehat J_{\alpha,g}$ avec la mesure et le noyau rend \eqref{eq:pdf2-ym-owner-stack-transfer-identity} compatible avec toute superposition. Sa limite définit
\begin{equation}\label{eq:pdf2-ym-owner-schwinger-functionals}
 S_n^{\rm YM}(A_1,t_1;\ldots;A_n,t_n)
 :=\lim_{\alpha}S_{\alpha,n}^{(\nu,g)}
 (A_{1,\alpha},t_1;\ldots;A_{n,\alpha},t_n).
\end{equation}

Les insertions de $\Phi$ s’obtiennent en dérivant en $t=0$ les unitaires
\eqref{eq:pdf2-ym-owner-curvature-weyl}. Puisque les $Z_{\alpha,c}$ transportés
sont centrés, indépendants et uniformément bornés, le lemme de Hoeffding donne,
niveau par niveau,
\begin{equation}\label{eq:pdf2-ym-owner-subgaussian-bound}
 \left\langle\Omega_\alpha,
 e^{t\Phi_\alpha(h)}\Omega_\alpha\right\rangle
 \le \exp\!\left(\frac9{16}t^2\|P_\alpha h\|_{L^2}^2\right).
\end{equation}
La borne passe à la limite par
\eqref{eq:pdf2-ym-owner-smeared-cauchy-derived} ; toutes les dérivées sont continues dans
la topologie de Schwartz et définissent des distributions tempérées.

La forme à huit faces se prolonge d’abord à toute pile finie. Si $F$ dépend des
$N$ couches positives d’un écran $\alpha$, la réflexion de
\eqref{eq:pdf2-ym-owner-markov-stack-law} et le conditionnement sur la couche centrale donnent
\begin{equation}\label{eq:pdf2-ym-owner-finite-stack-os}
 \begin{aligned}
 &\int \overline{F(\Theta_\nu y_{-N},\ldots,
                         \Theta_\nu y_{-1})}
          F(y_1,\ldots,y_N)\,d\Omega_{\alpha,N}^{(\nu,g)}(y)\\
 &\qquad=
 \int\nu_{\alpha,g}(dz)\left|
  \int K_{\alpha,g}(a_\alpha/2;z,dy_1)
       \prod_{j=1}^{N-1}K_{\alpha,g}(a_\alpha;y_j,dy_{j+1})
       F(y_1,\ldots,y_N)
 \right|^2\ge0.
 \end{aligned}
\end{equation}
Tout cylindre de demi-espace appartient à une pile finie ; par conséquent, \eqref{eq:pdf2-ym-owner-finite-stack-os} matérialise la forme OS avant l'adhérence. Dans le quotient OS, si $F$ possède des insertions à des temps positifs, soit $(\tau_sF)(A_1,t_1;\ldots;A_n,t_n):=
F(A_1,t_1+s;\ldots;A_n,t_n+s)$. On définit l'opérateur de translation par
\begin{equation}\label{eq:pdf2-ym-owner-os-transfer-operator}
 T_{\rm OS}^{(\nu)}(s)[F]:=[\tau_sF],
 \qquad s\ge0.
\end{equation}
Pour des cylindres $F,G$ de demi-espace, l’égalité euclidienne--Markov
\eqref{eq:pdf2-ym-owner-euclidean-markov-identification} et la propriété de
semi-groupe donnent générateur par générateur
\begin{equation}\label{eq:pdf2-ym-owner-os-markov-intertwining}
 \begin{aligned}
 \langle[F],T_{\rm OS}^{(\nu)}(s)[G]\rangle_{\rm OS}
 &=S_g^{\rm YM}((\Theta_\nu F)^*\tau_sG)\\
 &=\left\langle
    \mathcal V_{\infty,\nu}[F],
    e^{-sD_{\rm curv,g}^{\rm YM}}
    \mathcal V_{\infty,\nu}[G]
   \right\rangle .
 \end{aligned}
\end{equation}
L’identité montre à la fois que les classes nulles restent nulles et que,
par densité des cylindres,
\begin{equation}\label{eq:pdf2-ym-owner-os-generator-identification}
 \mathcal V_{\infty,\nu}T_{\rm OS}^{(\nu)}(s)
 \mathcal V_{\infty,\nu}^{-1}
 =e^{-sD_{\rm curv,g}^{\rm YM}}.
\end{equation}
C’est la flèche explicite mesure euclidienne $\to$ désintégration $\to$
transfert OS $\to$ hamiltonien.

\begin{proposition}[Paquet d’Osterwalder--Schrader continu]
\label{prop:pdf2-ym-owner-continuum-os-package}
Les fonctionnelles \eqref{eq:pdf2-ym-owner-schwinger-functionals} satisfont :
\begin{enumerate}
 \item normalisation, hermiticité et symétrie bosonique ;
 \item covariance euclidienne et régularité tempérée ;
 \item positivité par réflexion dans chacune des quatre directions,
 \begin{equation}\label{eq:pdf2-ym-owner-continuum-reflection-positive}
  \sum_{i,j}\overline{a_i}a_j\,
  S^{\rm YM}\!\left((\Theta_\nu F_i)^*F_j\right)\ge0,
  \qquad F_i\in\mathfrak A_{+,\nu};
 \end{equation}
 \item localité et propriété de cluster exponentielle : si $A,B$ sont locaux et centrés et si une
 translation de longueur $r$ sépare leurs supports, alors
 \begin{equation}\label{eq:pdf2-ym-owner-continuum-cluster}
  |S^{\rm YM}(A\,\tau_rB)|
  \le e^{-3\kappa_{\rm av}r}\,
       \|A\Omega_\infty\|\,\|B\Omega_\infty\|.
 \end{equation}
\end{enumerate}
La reconstruction OS de l'un quelconque des quatre demi-réseaux produit le même espace cyclique, le même vide et l'hamiltonien $D_{\rm curv,g}^{\rm YM}$. Le prolongement OS produit des champs de Wightman locaux et une représentation de Poincaré dont le spectre satisfait
\begin{equation}\label{eq:pdf2-ym-owner-wightman-spectrum}
 \operatorname{Spec}(P^0,\mathbf P)\subseteq\overline V_+,
 \qquad P^0=D_{\rm curv,g}^{\rm YM}\ge0.
\end{equation}
\end{proposition}

\begin{proof}
La normalisation et l’hermiticité proviennent de l’état ; la symétrie et la localité, de la
commutativité euclidienne à supports disjoints. La covariance et la continuité sont
\eqref{eq:pdf2-ym-owner-field-e4-covariance}--
\eqref{eq:pdf2-ym-owner-strong-e4} ; la régularité est
\eqref{eq:pdf2-ym-owner-subgaussian-bound}. Pour
$F_i$ cylindriques, tous vivent dans une superposition et une pile finies ;
\eqref{eq:pdf2-ym-owner-finite-stack-os} et la polarisation prouvent
\eqref{eq:pdf2-ym-owner-continuum-reflection-positive}. L’approximation en norme OS
et Cauchy--Schwarz l’étendent à $\mathfrak A_{+,\nu}$.

Après une rotation euclidienne, la séparation spatiale devient un temps positif.
L’identité pile--transfert
\eqref{eq:pdf2-ym-owner-stack-transfer-identity} et l’opérateur
\eqref{eq:pdf2-ym-owner-os-transfer-operator} écrivent la corrélation centrée sous la forme
\[
 \langle A^*\Omega_\infty,
 e^{-rD_{\rm curv,g}^{\rm YM}}B\Omega_\infty\rangle.
\]
La borne spectrale \eqref{eq:pdf2-ym-owner-curvature-sector-gap} donne \eqref{eq:pdf2-ym-owner-continuum-cluster}. Enfin, les connecteurs \eqref{eq:pdf2-ym-owner-os-connectors} forment un système dirigé ; leur réunion est dense d'après \eqref{eq:pdf2-ym-owner-smooth-core}. L'identité \eqref{eq:pdf2-ym-owner-common-os-hamiltonian} identifie les quatre adhérences et leurs générateurs à $D_{\rm curv,g}^{\rm YM}$. Les hypothèses déjà vérifiées de régularité, de covariance, de réflexion et de cluster sont les entrées du prolongement OS ; son théorème de reconstruction donne la localité de Wightman et \eqref{eq:pdf2-ym-owner-wightman-spectrum}.
\end{proof}

\subsection{Publication typée comme théorie de Yang--Mills}

L’identification ne repose pas seulement sur l’évaluation d’un clover sur une connexion donnée.
Soit $\mathcal R_G^{\rm loc}$ la $*$-algèbre de lecteurs formée par toutes les
fonctions de classe des holonomies d’incidence et de plaquette, leurs transports vers un
sommet de base et les polarisations du lecteur $\mathcal P^{F^2}$. Soit
\begin{equation}\label{eq:pdf2-ym-owner-curvature-fibre}
 \mathcal E_G:=\Lambda^2(\mathbb R^4)^*\otimes\mathfrak g
\end{equation}
avec le produit invariant. Dans chaque cellule, la polarisation de
$\mathcal P^{F^2}$ donne un gramien positif $G_{\alpha,c}^{F}$ sur
$\mathcal E_G$. On quotiente par son noyau et on applique sa racine positive ; le lecteur
normalisé résultant est noté $Y_{\alpha,c}(v)$ et vérifie
\begin{equation}\label{eq:pdf2-ym-owner-curvature-cell-gram}
 \langle Y_{\alpha,c}(v),Y_{\alpha,d}(w)\rangle
 =\delta_{cd}\langle v,w\rangle_{\mathcal E_G}.
\end{equation}
Si $d\subset c$, soit $T_{d\leftarrow c}$ le transport de repère déjà présent dans l'holonomie de surface. L'opérateur d'entrelacement de courbure complète est
\begin{equation}\label{eq:pdf2-ym-owner-curvature-intertwiner}
 I_{\alpha\beta}^{F}Y_{\alpha,c}(v)
 :=\sum_{d\subset c}\sqrt{\frac{|d|}{|c|}}\,
 Y_{\beta,d}\!\left(\operatorname{Ad}_{T_{d\leftarrow c}}v\right).
\end{equation}
L’invariance du produit de $\mathfrak g$, la disjonction des sous-cellules et la
composition des transports prouvent
\begin{equation}\label{eq:pdf2-ym-owner-curvature-intertwiner-identities}
 (I_{\alpha\beta}^{F})^*I_{\alpha\beta}^{F}=I,
 \qquad
 I_{\beta\gamma}^{F}I_{\alpha\beta}^{F}=I_{\alpha\gamma}^{F}.
\end{equation}
Pour $h\in C_c^\infty(\mathbb R^4;\mathcal E_G)$, on pose
\begin{equation}\label{eq:pdf2-ym-owner-f-alpha}
 \mathbf F_\alpha(h):=
 \sum_{c\in\mathcal C_\alpha^{\rm cell}}
 Y_{\alpha,c}\!\left(
   |c|^{-1/2}\int_c h(x)\,d^4x\right).
\end{equation}
Dans une superposition commune, le même calcul de
\eqref{eq:pdf2-ym-owner-smeared-cross-gram} donne maintenant
\begin{equation}\label{eq:pdf2-ym-owner-curvature-cross-gram}
 \left\langle I_{\alpha\gamma}^{F}\mathbf F_\alpha(h),
 I_{\beta\gamma}^{F}\mathbf F_\beta(k)\right\rangle
 =\langle P_\alpha h,P_\beta k\rangle_{L^2(\mathcal E_G)}.
\end{equation}
Par conséquent, $\mathbf F_\alpha(h)$ est de Cauchy par convergence forte des projections,
et non par une nouvelle hypothèse. Le même appariement de Haar et des innovations de
\eqref{eq:pdf2-ym-owner-publication-into-physical}, appliqué maintenant à la base
orthonormalisée de $G_{\alpha,c}^{F}$, réalise la colimite par une isométrie
locale et covariante de jauge dans
$L^2(\widehat X_\infty,\widehat\mu_{\infty,g}^{\rm YM})$. Sa limite définit la
distribution aléatoire
covariante de jauge
\begin{equation}\label{eq:pdf2-ym-owner-distributional-curvature}
\begin{aligned}
 \mathbf F_{\mu\nu}&\in
 \mathcal S'(\mathbb R^4;\mathfrak g)\ \widehat\otimes\
 L^2(\widehat X_\infty,\widehat\mu_{\infty,g}^{\rm YM}),\\
 U(a)\mathbf F(h)U(a)^*&=\mathbf F(a\!\cdot h),\\
 u\mathbf F(h)u^{-1}&=\mathbf F(\operatorname{Ad}_u h),
 \qquad a\in E(4),\ u\in G.
\end{aligned}
\end{equation}
Son canal de classe spectrale d'ordre trois est précisément $\Phi$. La relation de surface finie
\begin{equation}\label{eq:pdf2-ym-owner-surface-product}
 \operatorname{Hol}^{(\alpha)}_p
 =\prod_{p'\subset p}^{\longrightarrow}
  T_{p'}\operatorname{Hol}^{(\beta)}_{p'}T_{p'}^{-1}
\end{equation}
passe à la fermeture cofinale car tous ses produits croisés sont calculés dans une
superposition. En particulier, le lecteur clover aléatoire, et non seulement son évaluation dans une
connexion lisse, satisfait l’identité des petites boucles
\begin{equation}\label{eq:pdf2-ym-owner-holonomy-curvature-relation}
 \mathbf F(h)=L^2\!\mathord{-}\!\lim_{\alpha}
 a_\alpha^4\sum_x
 \left\langle(P_\alpha h)(x),F_{{\rm cl},\alpha}(x)\right\rangle,
 \qquad
 \delta_{\rm gauge}S_n^{\rm YM}=0,
 \qquad
 \operatorname{Alt}_{\lambda\mu\nu}D_\lambda\mathbf F_{\mu\nu}=0.
\end{equation}
La première égalité est \eqref{eq:pdf2-ym-owner-curvature-cross-gram} écrite dans le représentant clover ; la seconde est l'identité de Ward \eqref{eq:pdf2-ym-ward-finite} passée à la limite, et la troisième est l'annulation de la frontière d'une frontière dans le produit de surface. Aucune n'utilise le terme de masse.

Pour éviter de définir la cible par la conclusion que l’on veut obtenir, soit
$\mathbf{YM}_4(G,g)$ la classe des quintuplets locaux
$(\mathfrak A,\mu_g,K_g,\operatorname{Hol},\mathbf F)$ qui vérifient les
équations vérifiables suivantes : (i) la densité de Gibbs est
\eqref{eq:pdf2-ym-gibbs-measure} ; (ii) $K_g$ est markovien, réversible et satisfait
\eqref{eq:pdf2-ym-schwinger-dyson-finite}--\eqref{eq:pdf2-ym-ward-finite} ;
(iii) holonomie, clover et courbure satisfont
\eqref{eq:pdf2-ym-transport-surface-bianchi} et
\eqref{eq:pdf2-ym-owner-holonomy-curvature-relation} ; et (iv) leurs fonctionnelles sont la
même loi de piles de \eqref{eq:pdf2-ym-owner-stack-transfer-identity} et satisfont le
paquet OS\@. L’objet construit est, composante par composante,
\begin{equation}\label{eq:pdf2-ym-owner-qym-object}
\begin{aligned}
 \mathfrak Q_{\rm YM}^{(4)}(G,g):={}&
 \bigl(\{\mathfrak A_{\rm YM}(O)\}_O,
 S_g^{\rm YM},U,\Theta,\Omega_{\infty,g},\\
 &\widehat\mu_{\infty,g}^{\rm YM},
 \{S_{m,g},\mathscr S_{m,g}\}_m,K_g,
 \operatorname{Hol},\mathbf F,\\
 &\mathscr H_{\rm OS},D_{\rm curv,g}^{\rm YM}\bigr)
 \in\mathbf{YM}_4(G,g).
\end{aligned}
\end{equation}
L’application de publication n’omet aucune de ces composantes :
\begin{equation}\label{eq:pdf2-ym-owner-typed-publication-map}
 \begin{aligned}
 \mathcal P_{{\rm YM},g}:\mathfrak G_{\rm gau}(G)&\longrightarrow
 \mathbf{YM}_4(G,g),\\
 (\widehat\mu^0,\mathbb P^{\rm phys},D^0,\Theta,W,
   \operatorname{Hol},\mathcal P^{F^2})
 &\longmapsto
 (\widehat\mu_g^{\rm YM},S_g,\mathscr S_g,K_g,\mathbb P_g^{\rm phys},
  S_g^{\rm YM},\mathfrak A_{\rm YM},D_{\rm curv,g}^{\rm YM},
  \mathbf F,\operatorname{Hol}).
 \end{aligned}
\end{equation}
Les $S_n^{\rm YM}$ sont ainsi exactement les fonctions de Schwinger de la mesure dépendante de $g$, de la loi de Markov et de l'opérateur de transfert déjà construits.

Il reste à vérifier que $S_{m,g}$ est l’action de Yang--Mills, et non seulement une
fonction portant ce nom. Soit
$A\in C_c^3(\mathbb R^4;T^*\mathbb R^4\otimes\mathfrak g)$ une connexion lisse dans
une carte où le logarithme des petites boucles est unique. Pour la plaquette
$p_{\mu\nu}(x)$, le développement ordonné de Magnus donne
\begin{equation}\label{eq:pdf2-ym-owner-plaquette-magnus}
 \log\operatorname{Hol}_{p_{\mu\nu}(x)}(A)
 =a_m^2F_{A,\mu\nu}(x)
  +\frac{a_m^3}{2}(D_\mu+D_\nu)F_{A,\mu\nu}(x)
  +R_{m,\mu\nu}(x),
\end{equation}
avec $\|R_{m,\mu\nu}(x)\|\le C_G a_m^4
(\|D^2F_A\|_{\infty,p}+\|F_A\|_{\infty,p}^2)$. La moyenne des quatre orientations annule le terme impair et définit
\begin{equation}\label{eq:pdf2-ym-owner-clover}
 F_{{\rm cl},m}(x)
 :=\frac1{4a_m^2}\sum_{p\ni x}
  \operatorname{Ad}_{T_p}\log\operatorname{Hol}_p,
 \qquad
 \|F_{{\rm cl},m}(x)-F_A(x)\|
 \le C_G a_m^2(\|D^2F_A\|_\infty+\|F_A\|_\infty^2).
\end{equation}
La définition de la racine \eqref{eq:pdf2-ym-action-root} est précisément la polarisation de ces six composantes. En substituant \eqref{eq:pdf2-ym-owner-plaquette-magnus} dans chaque générateur de cette polarisation, on calcule
\begin{equation}\label{eq:pdf2-ym-owner-action-root-smooth-chart}
 \xi_{m,b}(A)
 =a_m^2\operatorname{vec}_{\mu<\nu}F_{A,\mu\nu}(x_b)
  +r_{m,b}(A),
 \qquad
 \|r_{m,b}(A)\|\le C_{G,A}a_m^4.
\end{equation}
Les termes croisés impairs s’annulent par les quatre orientations du clover ;
les termes pairs ne sont inclus qu’une seule fois par
\eqref{eq:pdf2-ym-action-factor-partition}. Par somme de Riemann et
Cauchy--Schwarz,
\begin{equation}\label{eq:pdf2-ym-owner-action-error}
 \left|S_{m,g}(A)-\frac1{4g^2}
  a_m^4\sum_x\|F_{{\rm cl},m}(x)\|^2\right|
 \le C_{G,A,g}a_m^2,
\end{equation}
et, en utilisant la borne de \eqref{eq:pdf2-ym-owner-clover},
\begin{equation}\label{eq:pdf2-ym-owner-f2}
 S_{m,g}(A)\longrightarrow
 S_{{\rm YM},g}(A):=\frac1{4g^2}
 \int_{\mathbb R^4}\|F_A(x)\|^2\,d^4x.
\end{equation}
Les équations de Schwinger--Dyson et de Ward ont déjà été calculées sur la mesure qui utilise
cette action. L’écart, en revanche, se déduit ensuite de la conjugaison unitaire de
\eqref{eq:pdf2-ym-g-dynamics} et du témoin
\eqref{eq:pdf2-ym-g-gap-before-limit} ; il n’a pas été utilisé pour choisir $g$, $S_{m,g}$ ni
$\widehat\mu_{m,g}^{\rm YM}$. Tel est le lien non circulaire entre action, loi et écart.

\begin{theorem}[Clôture de Yang--Mills pour tout groupe de Lie compact, connexe et simple]
\label{thm:pdf2-ym-cierre-global}
La structure complète $\mathfrak G_{\rm gau}$ construit une famille projective locale pour tout groupe de Lie compact, connexe et simple. L'application \eqref{eq:pdf2-ym-owner-typed-publication-map} publie une théorie quantique de Yang--Mills non triviale sur $\mathbb R^4$ : ses fonctionnelles satisfont l'ensemble continu des conditions OS et son prolongement satisfait les axiomes locaux de Wightman d'après la proposition~\ref{prop:pdf2-ym-owner-continuum-os-package} ; ses holonomies et sa courbure distributionnelle satisfont \eqref{eq:pdf2-ym-owner-holonomy-curvature-relation}, et les quatre reconstructions coinduisent le même hamiltonien $D_{\infty,g}^{\rm YM}$, avec un vide simple. Sa loi est $\widehat\mu_{\infty,g}^{\rm YM}$, construite par la famille compatible d'actions locales $\{S_{m,g},\mathscr S_{m,g}\}_m$, et satisfait les équations de Schwinger--Dyson et de Ward \eqref{eq:pdf2-ym-schwinger-dyson-finite}--\eqref{eq:pdf2-ym-ward-finite}. Le secteur cyclique de courbure \eqref{eq:pdf2-ym-owner-curvature-cyclic-sector} est réducteur, contient le témoin \eqref{eq:pdf2-ym-owner-wc-curvature-functional} et possède
\begin{equation}\label{eq:pdf2-ym-final-gap}
 \boxed{\Delta_{\rm YM}^{\rm HMT}=3\kappa_{\rm av}>0},
 \qquad
 m_{\rm gap}^{\rm HMT}
 =\frac{\hbar_{\rm int}}{c_{\rm int}}\Delta_{\rm YM}^{\rm HMT}>0.
\end{equation}
Le réseau local, l’action forte de $E(4)$, les distributions de Schwinger,
l’holonomie et $\mathbf F$ proviennent du même état complet. Le lecteur exact
$\mathcal P^{F^2}$ définit la densité de la loi ; le clover est sa carte lisse et
\eqref{eq:pdf2-ym-owner-f2} identifie son action classique. L’écart appartient au secteur collectif
de courbure invariant de jauge et n’attribue pas de masse élémentaire au gluon.
\end{theorem}

\begin{proof}
La projectivité cinématique est \eqref{eq:pdf2-ym-owner-projectivity} ; la racine
\eqref{eq:pdf2-ym-action-root}, le cocycle
\eqref{eq:pdf2-ym-action-cocycle} et la densité
\eqref{eq:pdf2-ym-gibbs-measure} construisent d’abord l’action et la mesure au
couplage $g$. Le transport facteur par facteur
\eqref{eq:pdf2-ym-gibbs-transport-factor}--
\eqref{eq:pdf2-ym-full-algebra-identities} prouve la projectivité, la préservation de
l’état, des produits et de $*$ ;
\eqref{eq:pdf2-ym-transport-surface-bianchi} établit clover et Bianchi. L’indice exhaustif
\eqref{eq:pdf2-ym-owner-exhaustive-index}, la forme
\eqref{eq:pdf2-ym-owner-full-form} et le circuit
\eqref{eq:pdf2-ym-owner-local-circuit} construisent le générateur commun. La réduction
physique et son noyau quotient sont
\eqref{eq:pdf2-ym-owner-physical-reduction}--
\eqref{eq:pdf2-ym-owner-D}. La conjugaison
\eqref{eq:pdf2-ym-g-dynamics} les transporte à la mesure dépendant de $g$ sans
altérer leur spectre. La décomposition de Hoeffding prouve le vide et la borne, et
\eqref{eq:pdf2-ym-owner-wc-invariance}--
\eqref{eq:pdf2-ym-owner-wc-eigenvalue} exhibent un vecteur physique qui atteint le
seuil. L’entrelacement nonadique transporte forme, vide et témoin à la limite, où
\eqref{eq:pdf2-ym-owner-limit-gap}--
\eqref{eq:pdf2-ym-owner-exact-gap} fixent l’écart exact. Les quatre
factorisations \eqref{eq:pdf2-ym-owner-four-os}, l’inclusion des images
\eqref{eq:pdf2-ym-owner-os-range-inclusion} et les connecteurs
\eqref{eq:pdf2-ym-owner-os-connectors} identifient leurs reconstructions à ce
même générateur. L’holonomie
\eqref{eq:pdf2-ym-owner-incidence-holonomy} place matériellement $W_c$ dans le secteur
de courbure, où \eqref{eq:pdf2-ym-owner-curvature-sector-gap} prouve l’écart
exact. Le système \eqref{eq:pdf2-ym-owner-block-intertwiner} contrôle les produits
croisés au moyen de \eqref{eq:pdf2-ym-owner-smeared-cross-gram} ;
\eqref{eq:pdf2-ym-full-publication-restriction}--
\eqref{eq:pdf2-ym-full-publication-products} élèvent ce premier chaos à l’isomorphisme
de toute l’algèbre. La convergence forte de $P_\alpha$ produit $\Phi$ ; les bornes
\eqref{eq:pdf2-ym-owner-holonomy-e4-estimate} et
\eqref{eq:pdf2-ym-owner-orbit-smearing-continuity} prouvent en outre la continuité
$E(4)$ des holonomies, des plaquettes et de clover. La désintégration euclidienne
\eqref{eq:pdf2-ym-owner-euclidean-markov-identification}, l’identité exacte
\eqref{eq:pdf2-ym-owner-stack-transfer-identity} et l’entrelacement
\eqref{eq:pdf2-ym-owner-os-markov-intertwining} relient la même mesure,
l’état de Schwinger, la pile de Markov, le transfert OS et le générateur de l’écart. La
proposition~\ref{prop:pdf2-ym-owner-continuum-os-package} prolonge les
quatre formes finies aux
demi-réseaux complets, prouve la régularité et la propriété de cluster et reconstruit l’hamiltonien. Pour
finir, \eqref{eq:pdf2-ym-owner-distributional-curvature}--
\eqref{eq:pdf2-ym-owner-typed-publication-map} identifient ces mêmes fonctionnelles
comme les fonctionnelles de Schwinger de la théorie de jauge locale ; les intégrations par parties
\eqref{eq:pdf2-ym-schwinger-dyson-finite}--\eqref{eq:pdf2-ym-ward-finite} prouvent
Schwinger--Dyson et Ward. Enfin,
\eqref{eq:pdf2-ym-owner-plaquette-magnus}--\eqref{eq:pdf2-ym-owner-f2} identifient
l’action classique avant d’appliquer l’équivalence spectrale. La chaîne n’utilise pas
l’écart pour choisir la mesure.
\end{proof}

\begin{falsifier}[Effondrement de la dynamique produit]
Si l’on remplace $D_m^{\rm prod}$ par $D_m^{\rm fib}$, tous les secteurs non
constants reçoivent la même valeur propre. Un vecteur de Hoeffding avec $|S|=2$ devrait
avoir l’énergie $6\kappa_{\rm av}$ par \eqref{eq:pdf2-ym-hoeffding}, mais le
semi-groupe simple lui attribuerait $3\kappa_{\rm av}$. L’identification des deux
générateurs est fausse.
\end{falsifier}

\begin{falsifier}[Confusion entre inclusion généalogique et moyenne de bloc]
Si l'on remplace $I_{\alpha\beta}^{\rm curv}$ par $\widehat J_{\alpha\beta}$ dans \eqref{eq:pdf2-ym-owner-smeared-cross-gram}, une coordonnée grossière persiste avec un coefficient égal à un tandis que le coefficient fin change selon le quotient des échelles. Pour $a_{m+1}=a_m/9$, le produit croisé ancestral porte un facteur $1/81$ et ne converge pas vers la norme diagonale. L'identité \eqref{eq:pdf2-ym-owner-smeared-cauchy-derived} échoue ; c'est pourquoi les deux applications restent séparées.
\end{falsifier}

\begin{falsifier}[Ablation de la fibre d’incidence]
Si l’on remplace $\widehat X_m$ par sa marginale de liens, on efface $q_c$ et, avec elle,
l’observable physique $W_c$. L’application d’oubli ne reflète pas la sous-algèbre de jauge complète ;
elle ne peut donc pas être utilisée pour nier la valeur propre de l’objet source. Objet
remplacé ; conclusion non transférable.
\end{falsifier}

\begin{falsifier}[Remplacement du lecteur exact par sa carte lisse]
Si l’on utilise directement le terme asymptotique de
\eqref{eq:pdf2-ym-owner-f2} comme densité sur un écran fini, on efface le reste
contrôlé de \eqref{eq:pdf2-ym-owner-action-error} et l’on perd le cocycle exact
\eqref{eq:pdf2-ym-action-cocycle}. La densité finie régissante est
$Z_{m,g}^{-1}e^{-\mathcal P_m^{F^2}/(4g^2)}$ ; le clover et la limite lisse
l’identifient à l’action classique, mais ne la remplacent pas.
\end{falsifier}

\paragraph{Provenance et statut.}
La séparation entre fibre simple et générateur produit, la mesure commune, la compression
de jauge, les quatre OS, la forme limite et le témoin physique sont
\texttt{RÉSULTAT\_RÉCUPÉRÉ} et \texttt{EXACT\_INTERNE} dans le domaine de jauge
régulier HMT enregistré. La désintégration à temps non uniformes, l'indice
exhaustif des espérances, le noyau quotient, la matrice $B^+$, l'action explicite
sur la colimite, le système cofinal de blocs, le paquet OS continu, la courbure
distributionnelle et l'holonomie d'incidence sont
\texttt{FORMALISATION\_NOUVELLE} : ils typent et satisfont les liens qui auparavant n'apparaissaient
que nominalement, sans introduire le gap comme donnée. Leur composition complète
établit le théorème précédent.
La racine positive du lecteur $F^2$, sa densité de Gibbs, le transport local compatible,
les équations de Schwinger--Dyson/Ward, l'isomorphisme algébrique complet, les bornes
de continuité des holonomies et l'identité pile--transfert sont
\texttt{FORMALISATION\_NOUVELLE}. Leur rôle est de satisfaire matériellement les quatre
ponts ; l'identification clover--action classique est une conséquence du développement
de Magnus et non une prémisse spectrale.
